\section{AIME II 2019}
        \begin{enumerate}
        	\item (\textit{AIME II 2019}) Two different points, $C$ and $D$, lie on the same side of line $AB$ so that $\triangle ABC$ and $\triangle BAD$ are congruent with $AB=9,BC=AD=10$, and $CA=DB=17$. The intersection of these two triangular regions has area $\frac{m}{n}$, where $m$ and $n$ are relatively prime positive integers. Find $m+n$.


 
 \textbf{Answer: }059



\textbf{Solution}


                    \begin{figure}[H]
                        \centering
                        \includegraphics[width=0.5\textwidth]{Solution/AIME_II_2019_1_Solution_Figure_1.png}
                    \end{figure}
                    - Diagram by Brendanb4321

 
Extend $AB$ to form a right triangle with legs $6$ and 

 


	\item (\textit{AIME II 2019}) Lily pads $1,2,3,\ldots$ lie in a row on a pond. A frog makes a sequence of jumps starting on pad $1$. From any pad $k$ the frog jumps to either pad $k+1$ or pad $k+2$ chosen randomly with probability $\frac{1}{2}$ and independently of other jumps. The probability that the frog visits pad $7$ is $\frac{p}{q}$, where $p$ and $q$ are relatively prime positive integers. Find $p+q$.


 
 \textbf{Answer: }107



\textbf{Solution (Probability States)}

Let $P_n$ be the probability the frog visits pad $7$ starting from pad $n$. Then $P_7 = 1$, $P_6 = \frac12$, and $P_n = \frac12(P_{n + 1} + P_{n + 2})$ for all integers $1 \leq n \leq 5$. Working our way down, we find
 \[P_5 = \frac{3}{4}\]
 \[P_4 = \frac{5}{8}\]
 \[P_3 = \frac{11}{16}\]
 \[P_2 = \frac{21}{32}\]
 \[P_1 = \frac{43}{64}\]$43 + 64 = \boxed{107}$.

 



\textbf{Solution 2 (Casework)}

Define a one jump to be a jump from $k$ to $k + 1$ and a two jump to be a jump from $k$ to $k + 2$. 

 Case 1: (6 one jumps) $\left (\frac{1}{2} \right)^6 = \frac{1}{64}$

 Case 2: (4 one jumps and 1 two jumps) $\binom{5}{1} \cdot \left(\frac{1}{2}\right)^5 = \frac{5}{32}$

 Case 3: (2 one jumps and 2 two jumps) $\binom{4}{2} \cdot \left(\frac{1}{2}\right)^4 = \frac{3}{8}$

 Case 4: (3 two jumps) $\left(\frac{1}{2}\right)^3 = \frac{1}{8}$

 Summing the probabilities gives us $\frac{43}{64}$ so the answer is $\boxed{107}$. 

 - pi\_is\_3.14

 



\textbf{Solution 3}

Let $P_n$ be the probability that the frog lands on lily pad $n$. The probability that the frog never lands on pad $n$ is $\frac{1}{2}P_{n-1}$, so $1-P_n=\frac{1}{2}P_{n-1}$. This rearranges to $P_n=1-\frac{1}{2}P_{n-1}$, and we know that $P_1=1$, so we can compute $P_7$.
 \begin{align*} P_1&=1\\ P_2&=1-\dfrac{1}{2} \cdot 1=\dfrac{1}{2}\\ P_3&=1-\dfrac{1}{2} \cdot \dfrac{1}{2}=\dfrac{3}{4}\\ P_4&=\dfrac{5}{8}\\ P_5&=\dfrac{11}{16}\\ P_6&=\dfrac{21}{32}\\ P_7&=\dfrac{43}{64}\\ \end{align*}
We calculate $P_7$ to be $\frac{43}{64}$, meaning that our answer is $\boxed{107}$.

 



\textbf{Solution 4}

For any point $n$, let the probability that the frog lands on lily pad $n$ be $P_n$. The frog can land at lily pad $n$ with either a double jump from lily pad $n-2$ or a single jump from lily pad $n-1$. Since the probability when the frog is at $n-2$ to make a double jump is $\frac{1}{2}$ and same for when it's at $n-1$, the recursion is just $P_n = \frac{P_{n-2}+P_{n-1}}{2}$. Using the fact that $P_1 = 1$, and $P_2 = \frac{1}{2}$, we find that $P_7 = \frac{43}{64}$. $43 + 64 = \boxed{107}$

 -bradleyguo

 



\textbf{Solution 5}

Let $a$ represent the number of two-step jumps (hopping over two lily pads) and $b$ represent the number of one-step jumps (hopping over one lily pad). We get the equation $2a + b = 6$, which has 4 integer solutions: $(3,0)$, $(2,2)$, $(1,4)$, and $(0,6)$. For $(2,2)$, there are $\frac{4!}{2!2!} = 6$ ways to permute the jumps, and for $(1,4)$ there are $\frac{5!}{1!4!} = 5$ ways. So we get:

 
 \[\frac{1}{2^{3+0}} + 6 \cdot \frac{1}{2^{2 + 2}} + 5 \cdot \frac{1}{2^{1 +4}} + \frac{1}{2^{0 + 6}} = \frac{43}{64}\]
$43 + 64 = \boxed{107}$




\textbf{Solution 6 (Rigorous Caseworks)}

Let $a$ be the number of 2-unit jumps and $b$ the number of 1-unit jumps. In this context, note that the total distance from 1 to 7 must be $7 - 1 = 6$, yielding the equation $2a + b = 6$ (i) in the process. On the other hand, we must establish a minimum and maximum limit for the jumps; thus, $a + b = m$ (ii), where $m \in \mathbb{Z}^+$. From there, we see:

 $\begin{cases}  2a + b = 6 \\  a + b = m  \end{cases} \Rightarrow \boxed{S = \{a = 6 - m, b = 2m - 6\}}$

 Since $a, b \in \mathbb{N}$, it follows that the interval for $m$ is $3 \le m \le 6$. Therefore, there are 4 cases to consider when the frog jumps. Let's illustrate them:

 
(i) $m = 3$: We will have $a = 3$ and $b = 0$, indicating that there is only one possible path for it to travel. Thus, the probability is $P_1 = \left(\frac{1}{2}\right)^3 = \frac{1}{8}$.

 (ii) $m = 4$: We will have $a = 2$ and $b = 2$, indicating that $\binom{4}{2} = 6$ possible paths for it to travel. Thus, the probability is $P_2 = \left(\frac{1}{2}\right)^4 \times 6 = \frac{6}{16}$.

 (iii) $m = 5$: We will have $a = 1$ and $b = 4$, indicating that $\binom{5}{1} = 5$ possible paths for it to travel. Thus, the probability is $P_3 = \left(\frac{1}{2}\right)^5 \times 5 = \frac{5}{32}$.

 (iv) $m = 6$: We will have $a = 0$ and $b = 6$, indicating that there is only one possible path for it to travel. Thus, the probability is $P_4 = \left(\frac{1}{2}\right)^6 = \frac{1}{64}$.

 Thus, we have:

 $\frac{p}{q} = P_1 + P_2 + P_3 + P_4 \implies \frac{p}{q} = \frac{1}{8} + \frac{5}{32} + \frac{6}{16} + \frac{1}{64} \implies \frac{p}{q} = \frac{43}{64}$

 Therefore, $p + q = 43 + 64 \therefore \boxed{\mathbf{p + q = 107}}$.

 



\textbf{Solution 7  (Recursion)}

We can solve this problem by looking at the probabilities of landing on each individual lily pad. Lily pad number one has a probability of 1, lily pad number 2 has a probability of 1/2. We can see from quick casework on lily pad number 3(yielding 3/4) that the probability for any lily pad N is $\frac{N_{n-2}}{2}+\frac{N_{n-1}}{2}$. Quickly working this out, we see that the probability is $\frac{43}{64}$, so the answer is $\boxed{107}$..

 


	\item (\textit{AIME II 2019}) Find the number of $7$-tuples of positive integers $(a,b,c,d,e,f,g)$ that satisfy the following system of equations: 

 \[abc=70\]

 \[cde=71\]

 \[efg=72.\]

 \textbf{Answer: }096



\textbf{Solution 1}

Note $c \mid 70, 71 \implies c \mid \text{gcd} (70, 71) = 1$ which implies $c=1$. Similarly, $e=1$, so $d=71$. Now, $\tau (70) = 8, \tau(72) = 12$ (where $\tau$ is the divisor function), so we have $8$ choices for $(a, b)$ and $12$ choices for $(f, g)$. This gives us a total of $8\cdot 12 = \boxed{096}$ solutions.

 



\textbf{Solution 2}

We know that any two consecutive numbers are coprime. Using this, we can figure out that $c=1$ and $e=1$. $d$ then has to be $71$. Now we have two equations left. $ab=70$ and $fg=72$. To solve these we just need to figure out all of the factors. Doing the prime factorization of $70$ and $72$, we find that they have $8$ and $12$ factors, respectively. The answer is $8 \times 12=\boxed{096}$




\textbf{Solution 3}

$c|70,71$ and $\gcd(70,71)=1$, so $c=1$. By similar logic, we find that $e=1$. $\newline$ Since $c=e=1$, the number of solutions is the number of factors of $70$ multiplied by the number of factors of $72$. $70=2^1 \times 5^1 \times7^1$, giving $(1+1)^3=8$ factors. $72=2^3 \times 3^2$, giving $(3+1)\times(2+1)=12$ factors. $\newline$ Thus, the final answer is $8\times12=\boxed{96}$. $\newline$ ~pandapower24

 


	\item (\textit{AIME II 2019}) A standard six-sided fair die is rolled four times. The probability that the product of all four numbers rolled is a perfect square is $\frac{m}{n}$, where $m$ and $n$ are relatively prime positive integers. Find $m+n$.


 
 \textbf{Answer: }187



\textbf{Solution 1}

Notice that, other than the number 5, the remaining numbers 1, 2, 3, 4, 6 are only divisible by 2 and/or 3. We can do some cases on the number of 5's rolled (note that there are $6^4 = 1296$ outcomes).

 Case 1 (easy): Four 5's are rolled. This has probability $\frac{1}{6^4}$ of occurring.

 Case 2: Two 5's are rolled.

 Case 3: No 5's are rolled.

 To find the number of outcomes for the latter two cases, we will use recursion. Consider a 5-sided die with faces numbered 1, 2, 3, 4, 6. For $n \ge 1$, let $a_n$ equal the number of outcomes after rolling the die $n$ times, with the property that the product is a square. Thus, $a_1 = 2$ as 1 and 4 are the only possibilities.

 To find $a_{n+1}$ given $a_n$ (where $n \ge 1$), we observe that if the first $n$ rolls multiply to a perfect square, then the last roll must be 1 or 4. This gives $2a_n$ outcomes. Otherwise, the first $n$ rolls do not multiply to a perfect square ($5^n - a_n$ outcomes). In this case, we claim that the last roll is uniquely determined (either 2, 3, or 6). If the product of the first $n$ rolls is $2^x 3^y$ where $x$ and $y$ are not both even, then we observe that if $x$ and $y$ are both odd, then the last roll must be 6; if only $x$ is odd, the last roll must be 2, and if only $y$ is odd, the last roll must be 3. Thus, we have $5^n - a_n$ outcomes in this case, and $a_{n+1} = 2a_n + (5^n - a_n) = 5^n + a_n$.

 Computing $a_2$, $a_3$, $a_4$ gives $a_2 = 7$, $a_3 = 32$, and $a_4 = 157$. Thus for Case 3, there are 157 outcomes. For case 2, we multiply $a_2$ by $\binom{4}{2} = 6$ to distribute the two 5's among four rolls. Thus the probability is

 
 \[\frac{1 + 6 \cdot 7 + 157}{6^4} = \frac{200}{6^4} = \frac{25}{162} \implies m+n = \boxed{187}\]

 -scrabbler94

 



\textbf{Solution 2}

We can solve this without finding the amount of cases. Notice when we roll a 1 or a 4, it does not affect whether or not the product is a square number. We have a 1/3 chance of rolling either a 1 or 4. We have a 2/3 chance of rolling a 2,3,5 or 6. Let's call rolling 1 or 4 rolling a dud (a perfect square). 

 Probability of rolling 4 duds: $\left(\frac{1}{3}\right)^4$

 Probability of rolling 3 duds: $4 * \left(\frac{1}{3}\right)^3 * \frac{2}{3}$

 Probability of rolling 2 duds: $6 * \left(\frac{1}{3}\right)^2 * \left(\frac{2}{3}\right)^2$

 Probability of rolling 1 dud: $4 * \frac{1}{3} * \left(\frac{2}{3}\right)^3$

 Probability of rolling 0 duds: $\left(\frac{2}{3}\right)^4$

 Now we will find the probability of a square product given we have rolled each amount of duds

 Probability of getting a square product given 4 duds: 1

 Probability of getting a square product given 3 duds: 0 (you will have 1 non-dud and that's never going to be square)

 Probability of getting a square product given 2 duds: $\frac{1}{4}$ (as long as our two non-duds are the same, our product will be square)

 Probability of getting a square product given 1 dud: $\frac{3!}{4^3}$ =  $\frac{3}{32}$ (the only way to have a square  product is rolling a 2,3 and 6. There are 3! ways of doing that and a total of $4^3$ ways to roll 3 non-duds).

 Probability of getting a square product given 0 duds: $\frac{40}{4^4}$= $\frac{5}{32}$ (We can have any two non-duds twice. For example, 2,2,5,5. There are $\binom{4}{2} = 6$ ways of choosing which two non-duds to use and $\binom{4}{2} = 6$ ways of choosing how to arrange those 4 numbers. That gives us 6*6=36 combinations. We can also have 2,2,2,2 or 3,3,3,3 or 5,5,5,5 or 6,6,6,6. This gives us a total of 40 combinations).

 We multiply each probability of rolling k duds with the probability of getting a square product given k duds and then sum all the values.
 \[\left(\frac{1}{3}\right)^4  * 1 + 4 * \left(\frac{1}{3}\right)^3 * \frac{2}{3} * 0 + 6 * \left(\frac{1}{3}\right)^2 * \left(\frac{2}{3}\right)^2 * \frac{1}{4} + 4 * \frac{1}{3} * \left(\frac{2}{3}\right)^3 * \frac{3}{32} + \left(\frac{2}{3}\right)^4 * \frac{5}{32} = \frac{25}{162}.\]
$25+162$ = $\boxed{187}$ 

 -dnaidu (silverlizard)

 



\textbf{Solution 3}

Note that rolling a 1/4 will not affect whether or not the product is a perfect square. This means that in order for the product to be a perfect square, all non 1/4 numbers rolled must come in pairs, with the only exception being the triplet 2,3, 6. Now we can do casework:

 If there are four 1/4's, then there are $2^4=16$ combinations. If there are three 1/4's, then there are 0 combinations, because the fourth number isn't a square. If there are two 1/4's, there are $2^2=4$ ways to choose the two 1/4's, 4 ways to choose the remaining pair of numbers, and $\frac{4!}{2!2!}=6$ ways to arrange, so there are $4\cdot 4\cdot 6=96$ combinations for this case. If there is one 1/4, then there are 2 ways to choose whether it is a 1 or 4, and the remaining three numbers must be 2, 3, and 6, so there are $4!$ ways to order, meaning there are $2\cdot 4!=48$ combinations for this case. Our final case is if there are no 1/4's, in which case we must have two pairs. If the two pairs are of different numbers, then there $\binom{4}{2}$ to choose the numbers and $\frac{4!}{2!2!}=6$ ways to arrange them, so $6\cdot 6=36$. If all four numbers are the same there are $4$ combinations, so there are $4+36=40$ combinations for this case.

 Hence there are $16+0+96+48+40=200$ combinations where the product of the dice is a perfect square, and there are $6^4=1296$ total combinations, so the desired probability is $\frac{200}{1296}=\frac{25}{162}$, yielding an answer of $25+162=\boxed{187}$.

 -Stormersyle

 



\textbf{Solution 4 (Casework)}

Another way to solve this problem is to do casework on all the perfect squares from $1^2$ to $36^2$, and how many ways they can be ordered$1^2$- $1,1,1,1$- $1$ way.$2^2$- $4,1,1,1$ or $2,2,1,1$- $\binom{4}{2}+4=10$ ways.$3^2$- $3,3,1,1$- $\binom{4}{2}=6$ ways.$4^2$- $4,4,1,1$, $2,2,2,2$, or $2,2,4,1$- $\binom{4}{2}+1+12=19$ ways.$5^2$- $5,5,1,1$- $\binom{4}{2}=6$ ways.$6^2$- $6,6,1,1$, $1,2,3,6$, $2,3,2,3$, $3,3,4,1$- $2*\binom{4}{2}+4!+12=48$ ways.$7^2$- Since there is a prime greater than 6 in its prime factorization there are $0$ ways.$8^2$- $4,4,4,1$ or $2,4,2,4$- $\binom{4}{2}+4=10$ ways.$9^2$- $3,3,3,3$- $1$ way.$10^2$- $2,2,5,5$ or $1,4,5,5$- $6+12=18$ ways.$11^2$- $0$ ways for the same reason as $7^2$.$12^2$- $6,6,2,2$, $4,4,3,3$, $2,3,4,6$, or $1,4,6,6$- $2*\binom{4}{2}+4!+12=48$ ways.$13^2$- $0$ ways.$14^2$- $0$ ways.$15^2$- $3,3,5,5$- $\binom{4}{2}=6$ ways.$16^2$- $4,4,4,4$- $1$ way.$17^2$- $0$ ways.$18^2$- $3,3,6,6$- $\binom{4}{2}=6$ ways.$19^2$- $0$ ways.$20^2$- $4,4,5,5$- $\binom{4}{2}=6$ ways.$21^2$- $0$ ways.$22^2$- $0$ ways.$23^2$- $0$ ways.$24^2$-$4,4,6,6$- $\binom{4}{2}=6$ ways.$25^2$- $5,5,5,5$- $1$ way.$26^2$- $0$ ways.$27^2$- $0$ ways.$28^2$- $0$ ways.$29^2$- $0$ ways.$30^2$- $5,5,6,6$- $\binom{4}{2}$ ways.$31^2$- $0$ ways.$32^2$- $0$ ways.$33^2$- $0$ ways.$34^2$- $0$ ways.$35^2$- $0$ ways.$36^2$- $6,6,6,6$- $1$ way.

 There are $6^4=1296$ ways that the dice can land. Summing up the ways, it is easy to see that there are $200$ ways. This results in a probability of $\frac{200}{1296}=\frac{25}{162}\implies\boxed{187}$-superninja2000

 



\textbf{Solution 5 (Recursion)}

We can do recursion on the number of rolls to find the number of ways we can get $4$ rolls to multiply to a square.

 After $n$ rolls, let us say that the product is $p = 2^a3^b5^c$.

 Define the following:

 $A_{n} =$ the number of ways to have a product after $n$ rolls where $a$ is odd, and $b$, $c$ are even

 $B_{n} =$ the number of ways to have a product after $n$ rolls where $b$ is odd, and $a$, $c$ are even

 $C_{n} =$ the number of ways to have a product after $n$ rolls where $c$ is odd, and $a$, $b$ are even

 $D_{n} =$ the number of ways to have a product after $n$ rolls where $c$ is even, and $a$, $b$ are odd

 $E_{n} =$ the number of ways to have a product after $n$ rolls where $b$ is even, and $a$, $c$ are odd

 $F_{n} =$ the number of ways to have a product after $n$ rolls where $a$ is even, and $b$, $c$ are odd

 $G_{n} =$ the number of ways to have a product after $n$ rolls where $a, b,$ and $c$ are all odd

 $S_{n} =$ the number of ways to have a product after $n$ rolls where $a, b,$ and $c$ are all even (square!)

 We have the following equations after considering the possible values of the nth roll:

 
 \[A_{n} = S_{n-1}+B_{n-1}+D_{n-1}+E_{n-1}+2A_{n-1}\]

 \[B_{n} = A_{n-1}+D_{n-1}+F_{n-1}+S_{n-1}+2B_{n-1}\]

 \[C_{n} = S_{n-1}+E_{n-1}+F_{n-1}+G_{n-1}+2C_{n-1}\]

 \[D_{n} = S_{n-1}+A_{n-1}+B_{n-1}+G_{n-1}+2D_{n-1}\]

 \[E_{n} = A_{n-1}+C_{n-1}+F_{n-1}+G_{n-1}+2E_{n-1}\]

 \[F_{n} = B_{n-1}+E_{n-1}+C_{n-1}+G_{n-1}+2F_{n-1}\]

 \[G_{n} = C_{n-1}+D_{n-1}+F_{n-1}+E_{n-1}+2G_{n-1}\]

 \[S_{n} = A_{n-1}+C_{n-1}+B_{n-1}+D_{n-1}+2S_{n-1}\]

 We have the following values after considering the possible values of the 1st roll:

 
 \[A_1 = B_1 = C_1 = D_1 = 1; E_1 = F_1 = G_1 = 0; S_1 = 2\]

 After applying recursion twice, we get:

 
 \[A_2 = B_2 = D_2 = 6, C_2 = 4, E_2 = F_2 = G_2 = 2, S_2 = 8\]

 \[A_3 = B_3 = D_3 = 34, C_3 = 22, E_3 = F_3 = G_3 = 18, S_3 = 38\]

 Finally, we have $S_4 = 200$, $\frac{m}{n} = \frac{200}{1296} = \frac{25}{162}$meaning our answer is $\boxed{187}$.

 



\textbf{Solution 6}

Consider all the distinct "fundamental" groups of integers from $1$ to $6$ whose product is a perfect square. A "fundamental" group is one that cannot be broken into two smaller groups that each have a perfect square product. For example, $\{2,2\}$ is a fundamental group, while $\{3,3,4\}$ is not, because it can be broken up into $\{3,3\}$ and $\{4\}$.

 $1$ and $4$ are already perfect squares, so they each form a "fundamental" group and cannot belong in any other group. Pairs of the other $4$ numbers ($\{2,2\}$,$\{3,3\}$, etc. ) form "fundamental" groups as well. The last "fundamental" group is $\{2,3,6\}$. It can be easily seen that no more groups exist.

 Thus, we have the "fundamental" groups $\{1\}$,$\{4\}$,$\{2,2\}$,$\{3,3\}$,$\{5,5\}$,$\{6,6\}$, and $\{2,3,6\}$.

 We now consider the ways to use these groups to form a sequence of $4$ numbers whose product is a perfect square. To form a set, we can simply select zero to two groups of size $2$ or $3$ and fill in any remaining spots with $1$s and $4$s. We can do this in one of $5$ ways: Using only $1$s and $4$s, using one group of size $2$, using one group of size $3$, using two different groups of size $2$, and using the same group of size $2$ twice.

 If we only use $1$s and $4$s, each of the $4$ slots can be filled with one of the $2$ numbers, so there are $2^4=16$ possibilities.

 If we use one group of size $2$, there are $4$ options for the group to use, $\binom{4}{2}$ ways to place the two numbers (since they are identical), and $2^2$ ways to fill in the remaining slots with $1$s and $4$s, so there are $4\cdot\binom{4}{2}\cdot2^2=96$ possibilities.

 If we use one group of size $3$, there is only $1$ option for the group to use, $4\cdot3\cdot2$ ways to place the three numbers (since they are distinct), and $2$ ways to fill in the remaining slot, so there are $4\cdot3\cdot2\cdot2=48$ possibilities.

 If we use two different groups of size $2$, there are $\binom{4}{2}$ options for the groups to use and $\binom{4}{2}$ ways to place the four numbers (since there are $2$ groups of identical numbers, and one group's placement uniquely determines the other's), so there are $\binom{4}{2}\cdot\binom{4}{2}=36$ possibilities.

 If we use the same group of size $2$ twice, there are $4$ options for the group to use and $1$ way to place the four numbers (since they are all identical), so there are $4=4$ possibilities.

 This gives us a total of $16+96+48+36+4=200$ possibilities, and since there are $6^4=1296$ total sequences that can be rolled, the probability is equal to $\frac{200}{1296}=\frac{25}{162}$, so the answer is $25+162=\boxed{187}$. ~\href{https://artofproblemsolving.com/wiki/index.php?title=User:Emerald\_block}{emerald\_block}

 



\textbf{Solution 7(Generating functions)}

Let's look at the prime factorization of some of these rolls:

 \begin{verbatim}
01 => 2\textasciicircum{}0*3\textasciicircum{}0*5\textasciicircum{}0 02 => 2\textasciicircum{}1*3\textasciicircum{}0*5\textasciicircum{}0 03 => 2\textasciicircum{}0*3\textasciicircum{}1*5\textasciicircum{}0 04 => 2\textasciicircum{}2*3\textasciicircum{}0*5\textasciicircum{}0 05 => 2\textasciicircum{}0*3\textasciicircum{}0*5\textasciicircum{}1 06 => 2\textasciicircum{}1*3\textasciicircum{}1*5\textasciicircum{}0
\end{verbatim}

Now, using multi-variable generating functions, we get:

 \begin{verbatim}
f(x,y,z)=1+x+y+x\textasciicircum{}2+z+xy            | |           \| |/            'v' f(x,y,z)=1+x+y+z+x\textasciicircum{}2+xy   =(for our purposes)       =2+x+y+z+xy
\end{verbatim}

Let's square that!

 \begin{verbatim}
4+2x+2y+2z+2xy+2x+x\textasciicircum{}2+xy+xz+x\textasciicircum{}2y+2y+xy+y\textasciicircum{}2+yz+xy\textasciicircum{}2+2z+xz+yz+z\textasciicircum{}2+xyz+2xy+x\textasciicircum{}2y+xy\textasciicircum{}2+xyz+x\textasciicircum{}2y\textasciicircum{}2 Combining like terms . . . 4+4x+4y+4z+x\textasciicircum{}2+y\textasciicircum{}2+z\textasciicircum{}2+6xy+2xz+2yz+2xyz+2x\textasciicircum{}2y+2xy\textasciicircum{}2+x\textasciicircum{}2y\textasciicircum{}2
\end{verbatim}

Since we only want the parity of each of the exponents, we can collapse it again.

 \begin{verbatim}
8+6x+6y+4z+6xy+2xz+2yz+2xyz
\end{verbatim}

Last simplification: factor out a factor of two and save it for later. 

 \begin{verbatim}
4+3x+3y+2z+3xy+xz+yz+xyz
\end{verbatim}

Let's take a better approach this time.

 \begin{verbatim}
\_\_\_ term :4*4+3*3+3*3+2*2+3*3+1*1+1*1+1*1=16+09+09+04+09+01+01+01=50 \_\_x term :4*3+3*4+3*3+2*1+3*3+1*2+1*1+1*1=12+12+09+02+09+02+01+01=48 \_\_y term :4*3+3*3+3*4+2*1+3*3+1*1+1*2+1*1=12+09+12+02+09+01+02+01=48 \_\_z term :4*2+3*1+3*1+2*4+3*1+1*3+1*3+1*3=08+03+03+08+03+03+03+03=34 \_xy term :4*3+3*3+3*3+2*1+3*4+1*1+1*1+1*2=12+09+09+02+12+01+01+02=48 \_xz term :4*1+3*2+3*1+2*3+3*1+1*4+1*3+1*3=04+06+03+06+03+04+03+03=32 \_yz term :4*1+3*1+3*2+2*3+3*1+1*3+1*4+1*3=04+03+06+06+03+03+04+03=32 xyz term :4*1+3*1+3*1+2*3+3*2+1*3+1*3+1+4=04+03+03+06+06+03+03+04=32  Result: 50+48x+48y+34z+48xy+32xz+32yz+32xyz
\end{verbatim}

I know we could use vectors and dot products to make it look neater but \textit{come on}. It \textbf{already} looks neat enough.Also, we didn't \textit{need} the other parts, but it just <i>looks \textbf{nicer}</i>.Now let's stick back the two that turned into a four.

 \begin{verbatim}
200+192x+192y+136z+192xy+128xz+128yz+128xyz
\end{verbatim}

We seek the constant term which is 200. 200/1296=100/648=50/324=25/162, 25+162=187

 



\textbf{Solution 8}

There are a total of $6^4=1296$ possible die rolls.

 
We use casework:

 
\textbf{Case 1}: All 4 numbers are the same.There are obviously $6$ ways.

 
\textbf{Case 2}: Sets of 2 different numbers.

 A set of two different numbers is basically $(x,x,y,y)$ . There are a total of  $\frac{4!}{2!\cdot 2!}=6$ ways to arrange the numbers. 

 By listing these cases, we quickly see a pattern:

 $(1,1,2,2)$
$(1,1,3,3)$
$...$
$(1,1,6,6)$
$(2,2,3,3)$
$...$
$(2,2,6,6)$
$...$
$(5,5,6,6)$

 There are a total of $5+4+3+2+1=15$ cases. Multiplying this by $6$ yields $15\cdot 6=90$ ways.

 


 \textbf{Case 3}: Sets of numbers in the form of $(x,x,1,4)$

 A special case must be made for the number $4$ because $4$ itself is a perfect square.

 $(1,1,1,4)$ $4$
$(2,2,1,4)$ $12$
$(3,3,1,4)$ $12$
$(4,4,1,4)$ $4$
$(5,5,1,4)$ $12$
$(6,6,1,4)$ $12$

 
Summing these up yields a total of $4+12+12+4+12+12=56$ ways.

 


 \textbf{Case 4}: Sets with all 4 numbers different

 Note that the sets

 $(1,2,3,6)$
$(2,3,4,6)$

 Multiply to perfect square. The total of these cases are $24+24=48$.

 


 
Adding all these cases together yields $6+90+56+48=200$ ways that the product of the values of the die can be a perfect square. 

 
Therefore the probability is

 $\frac{200}{1296}=\frac{25}{162}$
$m+n = 25+162 = \boxed{187}$

 -elchen3\_f

 



\textbf{Solution 9 (Clean Generating Functions and Roots of Unity Filter)}

Recall that a number is a perfect square iff every prime in its prime factorization has even multiplicity. We can use multi-variable generating functions to track the contribution of each prime (namely 2, 3, 5) in one roll. Let $f(a,b,c)$ be the generating function such that the coefficient of each term $a^xb^yc^z$ is the number of ways to get a product of $2^x3^y5^z$ when rolling 4 dice.  For each roll, we can either roll a 1 (basically contribute nothing), roll a 2 (contribute a factor of 2), roll a 3 (contribute a factor of 3), roll a 4 (contribute two factors of 2), roll a 5 (contribute a factor of 5), or roll a 6 (contribute a factor of 2 and a factor of 3), so the generating function for one roll is $1+a+b+a^2+c+ab$. Hence, $f(a,b,c)= (1+a+b+a^2+c+ab)^4$. We seek the sum of coefficients of the $a^{2p}b^{2q}c^{2r}$ terms $(p, q, r, \in \mathbb{Z}*)$. $\newline$

 To filter out the desired coefficients, apply Roots of Unity filter 3 times. Filter out the coefficients with even $a$ multiplicity using 
 \[\frac{f(1, b, c)+f(-1, b, c)}{2}=\frac{(3+2b+c)^4+(c+1)^4}{2}=g(b,c).\] Similarly, now filter out the coefficients with even $b$ multiplicity with 
 \[\frac{g(1, c)+g(-1, c)}{2}=\frac{(c+5)^4+3(c+1)^4}{4}=h(c).\] Apply this one more time to get 
 \[\frac{h(1)+h(-1)}{2}=\frac{6^4+3(2)^4+4^4}{8}.\] The requested probability is $\frac{6^4+3(2)^4+4^4}{8\cdot 6^4}=\frac{25}{162} \implies \boxed{187}.$



	\item (\textit{AIME II 2019}) Four ambassadors and one advisor for each of them are to be seated at a round table with $12$ chairs numbered in order $1$ to $12$. Each ambassador must sit in an even-numbered chair. Each advisor must sit in a chair adjacent to his or her ambassador. There are $N$ ways for the $8$ people to be seated at the table under these conditions. Find the remainder when $N$ is divided by $1000$.


 
 \textbf{Answer: }520



\textbf{Solution 0}

Suppose first that the chairs weren't numbered, and rotation doesn't matter. Then if we let the ambassador-advisor pairs act as a single person, we essentially want to pick $4$ seats out of $8$ seats, and fixing one pair gives $7 \cdot 6 \cdot 5$ ways. However, now we want to number the seats, in which there are $12$ possibilities (for some pair of people, we choose the order and then we choose which even spot the ambassador sits in, giving $2\cdot 6.$) However, the order of the other pairs is then uniquely determined, so this method works, giving us $7 \cdot 6 \cdot 5 \cdot 12 = 2520$ ways. We are asked to find the remainder when $\textit{N}$ is divided by $1000$, which is $\boxed{520}$.~Maximilian113 ~Minor edits by Hiphopfrog1

 



\textbf{Solution 1}

Let us first consider the $4$ ambassadors and the $6$ even-numbered chairs. If we consider only their relative positions, they can sit in one of $3$ distinct ways: Such that the $2$ empty even-numbered chairs are adjacent, such that they are separated by one occupied even-numbered chair, and such that they are opposite each other. For each way, there are $4!=24$ ways to assign the $4$ ambassadors to the $4$ selected seats.

 In the first way, there are $6$ distinct orientations. The $4$ advisors can be placed in any of the $5$ odd-numbered chairs adjacent to the ambassadors, and for each placement, there is exactly one way to assign them to the ambassadors. This means that there are $24\cdot6\cdot\binom{5}{4}=720$ total seating arrangements for this way.

 In the second way, there are $6$ distinct orientations. $3$ advisors can be placed in any of the $4$ chairs adjacent to the "chunk" of $3$ ambassadors, and $1$ advisor can be placed in either of the $2$ chairs adjacent to the "lonely" ambassador. Once again, for each placement, there is exactly one way to assign the advisors to the ambassadors. This means that there are $24\cdot6\cdot\binom{4}{3}\cdot\binom{2}{1}=1152$ total seating arrangements for this way.

 In the third way, there are $3$ distinct orientations. For both "chunks" of $2$ ambassadors, $2$ advisors can be placed in any of the $3$ chairs adjacent to them, and once again, there is exactly one way to assign them for each placement. This means that there are $24\cdot3\cdot\binom{3}{2}\cdot\binom{3}{2}=648$ total seating arrangements for this way.

 Totalling up the arrangements, there are $720+1152+648=2520$ total ways to seat the people, and the remainder when $2520$ is divided by $1000$ is $\boxed{520}$. ~\href{https://artofproblemsolving.com/wiki/index.php?title=User:Emerald\_block}{emerald\_block}

 



\textbf{Solution 2}

\begin{verbatim}
										
                    \begin{figure}[H]
                        \centering
                        \includegraphics[width=0.5\textwidth]{Solution/AIME_II_2019_5_Solution_Figure_1.png}
                    \end{figure}
                    
\end{verbatim}

In the diagram, the seats are numbered 1...12.  Rather than picking seats for each person, however, each ambassador/assistant team picks a gap between the seats (A...L) and the ambassador sits in the even seat while the assistant sits in the odd seat.  For example, if team 1 picks gap C then Ambassador 1 will sit in seat 2 while assistant 1 will sit in seat 3.  No two teams can pick adjacent gaps.  For example, if team 1 chooses gap C then team 2 cannot pick gaps B,C or D.  In the diagram, the teams have picked gaps C, F, H and J.  Note that the gap-gaps - distances between the chosen gaps - (in the diagram, 2, 1, 1, 4) must sum to 8.  So, to get the number of seatings, we:

 
            \begin{enumerate}
        \item Choose a gap for team 1 ($12$ options)
        \item Choose 3 other gaps around the table with positive gap-gaps. Note that because the sum of the distances between the gaps must be equal to $8,$ we create a bijection - to partition the 8 balls into 4 urns, with each urn having a positive number of balls.  There is the same as partitioning 4 with 4 non-negative integers, and using stars-and-bars, this is $\dbinom{4+4-1}{4-1}=\dbinom{7}{3}=35$
        \item Place the other 3 teams in the chosen gaps ($6$ permutations)
            \end{enumerate}
            So, the total is $12\cdot35\cdot6=2520$.

 Therefore, the remainder is $\boxed{520}$.

 



\textbf{Solution 3 (PIE)}

There are $6\cdot5\cdot4\cdot3\cdot2^4$ (placing the ambassadors and then choosing the locations of the advisors) total ways to let everyone sit. However, this may lead to advisors sitting in the same chair, leading to awkward situations. So we find how many ways this happens. There are $\dbinom{4}{2} = 6$ ways to choose which advisors end up sitting together, times $6\cdot2=12$ ways to find neighboring even seats (order matters), and for the other two ambassadors, there are $\dbinom{4}{2} \cdot 2 = 12$ ways to sit and 4 ways for their advisors to sit. This results in $6\cdot12\cdot12\cdot4=3456$ ways for this to happen. However, we overcounted the case when all four of the advisors run out of room to sit, where there are $\frac{6\cdot3}{2!}=9$ to choose the places the ambassadors sit and 4! ways to arrange the ambassadors. $9\cdot4!=216$ ways to happen. So the answer is $5760 - 3456 + 216 = 2520$, which has a remainder of $\boxed{520}$ when divided by 1000.

 



\textbf{Solution 4}

We see that for every 2 adjacent spots on the table, there is exactly one way for an ambassador and his or her partner to sit. There are 2 cases:

 
            \begin{enumerate}
        \item There is an ambassador at the 12th chair and his or her partner at the 1st chair
        \item There is no pair that has their chairs numbered as 12 and 1
            \end{enumerate}
            For the first case, there are 3 pairs and 4 empty spaces, which results in $\frac{7!}{3!4!} = 35$ ways to arrange the pairs.

 For the second case, there are 4 pairs and 4 empty spaces, which results in $\frac{8!}{4!4!} = 70$ ways to arrange the pairs.

 Finally, we have to assign ambassadors and their advisors to the pairs, which has $4!=24$ ways, so $N = (35+70)\cdot 24 = 2520$, or $\boxed{520}$~PenguinMoosey

 



\textbf{Solution 5 (Lots of diagrams)}


                    \begin{figure}[H]
                        \centering
                        \includegraphics[width=0.5\textwidth]{Solution/AIME_II_2019_5_Solution_Figure_2.png}
                    \end{figure}
                    Let's begin with a simpler problem. We start with n ambassador chairs in a row and n+1 advisor chairs between them. As before, each advisor has to sit next to his/her ambassador. In the diagram, a black spot represents an ambassador, a blue ring represents an empty advisor seat, and a black dot on top of a blue dot represents an advisor. There are exactly (n+1)! solutions in this case as represented by the diagram. First, let's seat the ambassadors in their positions (n! ways). In the diagram, the first solution is that the advisors all sit an the left of their respective ambassadors. Then, we can slide the rightmost advisor over his/her ambassador to make a new solution. Then, we can do the same with the other advisor next to the empty seat. We can continue this process until all of the advisors sit an the right of their respective ambassadors. In total, there are n!(n+1) = (n+1)! ways. 

 
                    \begin{figure}[H]
                        \centering
                        \includegraphics[width=0.5\textwidth]{Solution/AIME_II_2019_5_Solution_Figure_3.png}
                    \end{figure}
                    Now, we are ready to tackle the original problem. If we forget about the numbers and the uniqueness of people, we will have the three configurations shown above. In the first configuration, there is a single red blob of four ambassadors. Let's stretch them out onto a line. Now, it's the same situation as the simpler problem and there are (4+1)! = 5! = 120 ways to seat the people. However, we need to multiply that by 6 because there are 6 orientations. Therefore, we have 720 outcomes for this case. In the second configuration, there is a orange blob and a green blob, each of the case n=2, resulting in (2+1)! * (2+1)! which equals 36. Multiplying the number of orientations, 3, we arrive at 108. Wait! We need to multiply by 6 to account for the many ways that we can choose 2 of the ambassadors to be in the orange blob. Therefore, we have 648 outcomes. Last but not least, we have the final configuration which is a blob of 1 and a blob of 3. Applying the formula results in (1+1)! * (3+1)! which evaluates to 48. Then, we multiply by 4, the number of ways to choose the blobs, and 6, the number of orientations. Our result is 1152. Finally, adding the numbers, we get 2520 and taking mod 1000, we get 520.

 


	\item (\textit{AIME II 2019}) In a Martian civilization, all logarithms whose bases are not specified are assumed to be base $b$, for some fixed $b\ge2$. A Martian student writes down

 \[3\log(\sqrt{x}\log x)=56\]

 \[\log_{\log x}(x)=54\]
and finds that this system of equations has a single real number solution $x>1$. Find $b$.


 
 \textbf{Answer: }216



\textbf{Solution 1}

Using change of base on the second equation to base b,
 \[\frac{\log x}{\log \log_{b}{x}}=54\]
 \[\log x = 54 \cdot \log\log_{b}{x}\]
 \[x = (\log_{b}{x})^{54}\]Note by dolphin7 - you could also just rewrite the second equation in exponent form.

 Substituting this into the $\sqrt x$ of the first equation,
 \[3\log_{b}{((\log_{b}{x})^{27}\log_{b}{x})} = 56\]
 \[3\log_{b}{(\log_{b}{x})^{28}} = 56\]
 \[\log_{b}{(\log_{b}{x})^{84}} = 56\]

 We can manipulate this equation to be able to substitute $x = (\log_{b}{x})^{54}$ a couple more times:
 \[\log_{b}{(\log_{b}{x})^{54}} = 56 \cdot \frac{54}{84}\]
 \[\log_{b}{x} = 36\]
 \[(\log_{b}{x})^{54} = 36^{54}\]
 \[x = 6^{108}\]

 However, since we found that $\log_{b}{x} = 36$, $x$ is also equal to $b^{36}$. Equating these,
 \[b^{36} = 6^{108}\]
 \[b = 6^3 = \boxed{216}\]




\textbf{Solution 2}

We start by simplifying the first equation to 
 \[3\log_{b}{(\sqrt{x}\log x)}=\log_{b}{(x^{\frac{3}{2}}\log^3x)}=56\]
 \[x^\frac{3}{2}\cdot \log_b^3x=b^{56}\]Next, we simplify the second equation to 
 \[\log_{\log(x)}(x)=\frac{\log_b(x)}{\log_b(\log_b(x))}=54\]
 \[\log_bx=54\log_b(\log_b(x))=\log_b(\log_b^{54}(x))\]
 \[x=\log_b^{54}x\]Substituting this into the first equation gives 
 \[\log_b^{54\cdot \frac{3}{2}}(x)\cdot \log_b^3x=\log_b^{84}x=b^{56}\]
 \[x=b^{b^{\frac{56}{84}}}=b^{b^{\frac{2}{3}}}\]Plugging this into $x=\log_b^{54}x$ gives 
 \[b^{b^{\frac{2}{3}}}=\log_b^{54}(b^{b^\frac{2}{3}})=b^{\frac{2}{3}\cdot 54}=b^{36}\]
 \[b^{\frac{2}{3}}=36\]
 \[b=36^{\frac{3}{2}}=6^3=\boxed{216}\]-ktong

 



\textbf{Solution 3}

Apply change of base to 
 \[\log_{\log x}(x)=54\] to yield: 
 \[\frac{\log_b(x)}{\log_b(\log_b(x))}=54\]which can be rearranged as:  
 \[\frac{\log_b(x)}{54}=\log_b(\log_b(x))\]Apply log properties to 
 \[3\log(\sqrt{x}\log x)=56\] to yield: 
 \[3(\frac{1}{2}\log_b(x)+\log_b(\log_b(x)))=56\Rightarrow\frac{1}{2}\log_b(x)+\log_b(\log_b(x))=\frac{56}{3}\]Substituting 
 \[\frac{\log_b(x)}{54}=\log_b(\log_b(x))\] into the equation $\frac{1}{2}\log_b(x)+\log_b(\log_b(x))=\frac{56}{3}$ yields: 
 \[\frac{1}{2}\log_b(x)+\frac{\log_b(x)}{54}=\frac{28\log_b(x)}{54}=\frac{56}{3}\]So 
 \[\log_b(x)=36.\]Substituting this back in to 
 \[\frac{\log_b(x)}{54}=\log_b(\log_b(x))\] yields
 \[\frac{36}{54}=\log_b(36).\]So,
 \[b^{\frac{2}{3}}=36\Rightarrow \boxed{b=216}\]

 -Ghazt2002

 



\textbf{Solution 4}

1st equation: 
 \[\log (\sqrt{x}\log x)=\frac{56}{3}\] 
 \[\log(\sqrt x)+\log(\log x)=\frac{1}{2}\log x+\log(\log x)=\frac{56}{3}\] 2nd equation: 
 \[x=(\log x)^{54}\] So now substitute $\log x=a$ and $x=b^a$: 
 \[b^a=a^{54}\] 
 \[b=a^{\frac{54}{a}}\] We also have that 
 \[\frac{1}{2}a+\log_{a^\frac{54}{a}} a=\frac{56}{3}\] 
 \[\frac{1}{2}a+\frac{1}{54}a=\frac{56}{3}\] This means that $\frac{14}{27}a=\frac{56}{3}$, so 
 \[a=36\] 
 \[b=36^{\frac{54}{36}}=36^\frac{3}{2}=\boxed{216}\].

 -Stormersyle

 



\textbf{Solution 5 (Substitution)}

Let $y = \log _{b} x$Then we have
 \[3\log _{b} (y\sqrt{x}) = 56\]
 \[\log _{y} x = 54\]which gives
 \[y^{54} = x\]Plugging this in gives
 \[3\log _{b} (y \cdot y^{27}) = 3\log _{b} y^{28} = 56\]which gives
 \[\log _{b} y = \dfrac{2}{3}\]so
 \[b^{2/3} = y\]By substitution we have
 \[b^{36} = x\]which gives
 \[y = \log _{b} x = 36\]Plugging in again we get
 \[b = 36^{3/2} = \fbox{216}\]

 --Hi3142

 



\textbf{Solution 6 (Also Substitution)}

This system of equations looks complicated to work with, so we let $a=\log_bx$ to make it easier for us to read. 

 Now, the first equation becomes $3\log(\sqrt x \cdot a) = 56 \implies \log(\sqrt{x}\cdot a)=\frac{56}3$. 

 The second equation, $\log_{\log(x)}(x)=54$ gives us $\underline{a^{54} = x}$.

 Let's plug this back into the first equation to see what we get: $\log_b(\sqrt{a^{54}}\cdot a)=\frac{56}3$, and simplifying, $\log_b(a^{27}\cdot a^1)=\log_b(a^{28})=\frac{56}{3}$, so $b^{\frac{56}3}=a^{28}\implies \underline{b^{\frac 23}=a}$.

 Combining this new finding with what we had above $a^{54} = (b^{\frac 23})^{54} = x\implies \mathbf{b^{36} =x}$. 

 Now that we've expressed one variable in terms of the other, we can plug this into either equation, say equation 1. Then we get $\log_b(\sqrt{b^{36}}\cdot\log_b(b^{36})=\frac{56}3\implies$$\log_b(b^{18}\cdot 36)=\frac{56}3\implies b^{\frac{56}3}=b^{18}\cdot 36$. 

 Finally, that gives us that $\frac{b^{\frac{56}3}}{b^{18}}=36\implies b^{\frac{56}{3}-18}=b^{\frac{56}{3}-\frac{54}{3}}=b^{\frac 23}=36\implies b=36^{\frac 32}=6^3$. Thus, $b=\boxed{216}$.

 



\textbf{Solution 7 (Easy System of Equations)}

Using change of base on the second equation, we have

 
 \[\frac{\log_{b} x}{\log_{b} \log_{b} x} = 54\]

 Using log rules on the first equation, we have

 
 \[\frac{3}{2} \log_{b} x + 3 \log_{b} \log_{b} x = 56\]

 We notice that $\log_{b} x$ and $\log_{b} \log_{b} x$ are in both equations. Thus, we set $m = \log_{b} x$ and $n = \log_{b} \log_{b} x$ and we have

 
 \[\frac{3}{2} m + 3n = 56\]
 \[\frac{m}{n} = 54\]

 Solving this yields $m = 36$, $n = \frac{2}{3}$.

 Now, $n = \log_{b} \log_{b} x = \log_{b} m = \log_{b} 36$, so we have $\log_{b} 36 = \frac{2}{3}$. Solving this yields $b = \boxed{216}$.

 



\textbf{Solution 8 (Definition of Logarithm)}

The second equation implies that
 \[\log_{\log_b x} x=54\implies (\log_b x)^{54}=x \implies \log_b x=x^{\frac{1}{54}}\]The first equation implies that
 \[3\log_b(\sqrt{x} \log_b x)=56 \implies b^{\frac{56}{3}}=\sqrt{x} \log_b x\]Substituting the first result into the second gives us 
 \[b^{\frac{56}{3}}=x^{\frac{1}{2}}\cdot x^{\frac{1}{54}} \implies b=x^{\frac{1}{36}}.\]Because $b^{36}=x,$ $\log_b x=36$ by the definition of a logarithm.Substituting this into the second equation,
 \[\log_{36} x=54 \implies x=36^{54}.\]Finally, 
 \[b=(36^{54})^{\frac{1}{36}}=36^{54\cdot\frac{1}{36}}=6^{2*\frac{3}{2}}=6^3=\boxed{216}.\]



	\item (\textit{AIME II 2019}) Triangle $ABC$ has side lengths $AB=120,BC=220$, and $AC=180$. Lines $\ell_A,\ell_B$, and $\ell_C$ are drawn parallel to $\overline{BC},\overline{AC}$, and $\overline{AB}$, respectively, such that the intersections of $\ell_A,\ell_B$, and $\ell_C$ with the interior of $\triangle ABC$ are segments of lengths $55,45$, and $15$, respectively. Find the perimeter of the triangle whose sides lie on lines $\ell_A,\ell_B$, and $\ell_C$.


 
 \textbf{Answer: }715



\textbf{Solution 1}

Let the points of intersection of 

 


	\item (\textit{AIME II 2019}) The polynomial $f(z)=az^{2018}+bz^{2017}+cz^{2016}$ has real coefficients not exceeding $2019,$ and $f\left(\frac{1+\sqrt3i}{2}\right)=2015+2019\sqrt3i$. Find the remainder when $f(1)$ is divided by $1000$.


 
 \textbf{Answer: }053



\textbf{Solution 1}

We have $\frac{1+\sqrt{3}i}{2} = \omega$ where $\omega = e^{\frac{i\pi}{3}}$ is a primitive 6th root of unity. Then we have

 
 \begin{align*} f(\omega) &= a\omega^{2018} + b\omega^{2017} + c\omega^{2016}\\ &= a\omega^2 + b\omega + c \end{align*}


 We wish to find $f(1) = a+b+c$. We first look at the real parts. As $\text{Re}(\omega^2) = -\frac{1}{2}$ and $\text{Re}(\omega) = \frac{1}{2}$, we have $-\frac{1}{2}a + \frac{1}{2}b + c = 2015 \implies -a+b + 2c = 4030$. Looking at imaginary parts, we have $\text{Im}(\omega^2) = \text{Im}(\omega) = \frac{\sqrt{3}}{2}$, so $\frac{\sqrt{3}}{2}(a+b) = 2019\sqrt{3} \implies a+b = 4038$. As $a$ and $b$ do not exceed 2019, we must have $a = 2019$ and $b = 2019$. Then $c = \frac{4030}{2} = 2015$, so $f(1) = 4038 + 2015 = 6053 \implies f(1) \pmod{1000} = \boxed{053}$.

 -scrabbler94

 



\textbf{Solution 2}

Denote $\frac{1+\sqrt{3}i}{2}$ with $\omega$.

 By using the quadratic formula ($\frac{-b\pm\sqrt{b^2-4ac}}{2a}$) in reverse, we can find that $\omega$ is the solution to a quadratic equation of the form $ax^2+bx+c=0$ such that $2a=2$, $-b=1$, and $b^2-4ac=-3$. This clearly solves to $a=1$, $b=-1$, and $c=1$, so $\omega$ solves $x^2-x+1=0$.

 Multiplying $x^2-x+1=0$ by $x+1$ on both sides yields $x^3+1=0$. Muliplying this by $x^3-1$ on both sides yields $x^6-1=0$, or $x^6=1$. This means that $\omega^6=1$.

 We can use this to simplify the equation $a\omega^{2018}+b\omega^{2017}+c\omega^{2016}=f(\omega)=2015+2019\sqrt{3}i$ to $a\omega^2+b\omega+c=2015+2019\sqrt{3}i.$

 As in Solution 1, we use the values $\omega=\frac{1}{2}+\frac{\sqrt{3}}{2}i$ and $\omega^2=-\frac{1}{2}+\frac{\sqrt{3}}{2}i$ to find that $-\frac{1}{2}a+\frac{1}{2}b+c=2015$ and $\frac{\sqrt{3}}{2}a+\frac{\sqrt{3}}{2}b=2019\sqrt{3} \implies a+b=4038.$ Since neither $a$ nor $b$ can exceed $2019$, they must both be equal to $2019$. Since $a$ and $b$ are equal, they cancel out in the first equation, resulting in $c=2015$.

 Therefore, $f(1)=a+b+c=2019+2019+2015=6053$, and $6053\bmod1000=\boxed{053}$. ~\href{https://artofproblemsolving.com/wiki/index.php?title=User:Emerald\_block}{emerald\_block}

 


 



\textbf{Solution 3}

Calculate the first few powers of $\frac{1+\sqrt{3}i}{2}$.

 $(\frac{1+\sqrt{3}i}{2})^1=\frac{1+\sqrt{3}i}{2}$
$(\frac{1+\sqrt{3}i}{2})^2=\frac{-1+\sqrt{3}i}{2}$
$(\frac{1+\sqrt{3}i}{2})^3=-1$
$(\frac{1+\sqrt{3}i}{2})^4=\frac{-1-\sqrt{3}i}{2}$
$(\frac{1+\sqrt{3}i}{2})^5=\frac{1-\sqrt{3}i}{2}$
$(\frac{1+\sqrt{3}i}{2})^6=1$

 We figure that the power of $\frac{1+\sqrt{3}i}{2}$ repeats in a cycle 6.

 $f(\frac{1+\sqrt{3}i}{2})=(a(\frac{1+\sqrt{3}i}{2})^2+b(\frac{1+\sqrt{3}i}{2})+c)(\frac{1+\sqrt{3}i}{2})^{2016}$

 Since 2016 is a multiple of 6, $(\frac{1+\sqrt{3}i}{2})^{2016}=1$
$f(\frac{1+\sqrt{3}i}{2})=a(\frac{-1+\sqrt{3}i}{2})+b(\frac{1+\sqrt{3}i}{2})+c$
$f(\frac{1+\sqrt{3}i}{2})=(-\frac{1}{2}a+\frac{1}{2}b+c)+(\frac{\sqrt{3}i}{2}a+\frac{\sqrt{3}i}{2}b)=2015+2019\sqrt{3}i$

 Therefore, $-\frac{1}{2}a+\frac{1}{2}b+c=2015$ and $\frac{\sqrt{3}i}{2}a+\frac{\sqrt{3}i}{2}b=2019\sqrt{3}i$

 Using the first equation, we can get that $-a+b+2c=4030$, and using the second equation, we can get that $a+b=4038$.

 Since all coefficients are less than or equal to $2019$, $a=b=2019$.

 Therefore, $2c=4030$ and $c=2015$.

 $f(1)=a+b+c=2019+2019+2015=6053$, and the remainder when it divides $1000$ is $\boxed{053}.$



	\item (\textit{AIME II 2019}) Call a positive integer $n$ $k$-\textit{pretty} if $n$ has exactly $k$ positive divisors and $n$ is divisible by $k$. For example, $18$ is $6$-pretty. Let $S$ be the sum of the positive integers less than $2019$ that are $20$-pretty. Find $\frac{S}{20}$.


 
 \textbf{Answer: }472



\textbf{Solution 1}

Every 20-pretty integer can be written in form $n = 2^a 5^b k$, where $a \ge 2$, $b \ge 1$, $\gcd(k,10) = 1$, and $d(n) = 20$, where $d(n)$ is the number of divisors of $n$. Thus, we have $20 = (a+1)(b+1)d(k)$, using the fact that the divisor function is multiplicative. As $(a+1)(b+1)$ must be a divisor of 20, there are not many cases to check.

 If $a+1 = 4$, then $b+1 = 5$. But this leads to no solutions, as $(a,b) = (3,4)$ gives $2^3 5^4 > 2019$.

 If $a+1 = 5$, then $b+1 = 2$ or $4$. The first case gives $n = 2^4 \cdot 5^1 \cdot p$ where $p$ is a prime other than 2 or 5. Thus we have $80p < 2019 \implies p = 3, 7, 11, 13, 17, 19, 23$. The sum of all such $n$ is $80(3+7+11+13+17+19+23) = 7440$. In the second case $b+1 = 4$ and $d(k) = 1$, and there is one solution $n = 2^4 \cdot 5^3 = 2000$.

 If $a+1 = 10$, then $b+1 = 2$, but this gives $2^9 \cdot 5^1 > 2019$. No other values for $a+1$ work.

 Then we have $\frac{S}{20} = \frac{80(3+7+11+13+17+19+23) + 2000}{20} = 372 + 100 = \boxed{472}$.

 -scrabbler94

 



\textbf{Solution 2}

For $n$ to have exactly $20$ positive divisors, $n$ can only take on certain prime factorization forms: namely, $p^{19}, p^9q, p^4q^3, p^4qr$ where $p,q,r$ are primes. No number that is a multiple of $20$ can be expressed in the first form because 20 has \textit{two} primes in its prime factorization, while the first form has only \textit{one}, and the only integer divisible by $20$ that has the second form is $2^{9}5$, which is 2560, greater than $2019$. 

 For the third form, the only $20$-pretty numbers are $2^45^3=2000$ and $2^35^4=5000$, and only $2000$ is small enough. 

 For the fourth form, any number of the form $2^45r$ where $r$ is a prime other than $2$ or $5$ will satisfy the $20$-pretty requirement. Since $n=80r<2019$, $r\le 25$. Therefore, $r$ can take on $3, 7, 11, 13, 17, 19,$ or $23$. 

 Thus, $\frac{S}{20}=\frac{2000+80(3+7+11+...+23)}{20}=100+4(3+7+11+...+23)=\boxed{472}$.

 Rephrased for clarity by Afly

 



\textbf{Solution 3}

We have the number being $20k = 2^2 \cdot 5 \cdot k$. We will consider some cases.

 Case 1: $k$ doesn't contain $2$ nor $5$ in its prime factorization

 Then let $k$ have $p$ divisors and so we have

 $3 \cdot 2 \cdot p = 6p$ as the total number of divisors of $20k$. Thus must be equal to $20$ which implies $p = \frac{10}{3}$ a contradiction. So $k$ must contain either $2$ or $5$ or both in its prime factorization.

 Case 2: $k$ contains $2$ but not $5$ in its prime factorization

 Then the number is $2^{2 + a} \cdot 5 \cdot \text{something}$ assuming $k$ has $a$ factors of $2$. That something is the rest of $k$. Let that something have $x$ factors. Thus we have

 $(3 + a)(2)(x) = 20 \implies x(3 + a) = 10$

 If $x = 1$ then $a = 7$ which implies the number is $2^9 \cdot 5$ which is too big.

 If $x = 2$ then $a = 2$ implying the number must be $2^4 \cdot 5 \cdot p$ since if $x = 2$ the rest of the something of $k$ must have $2$ divisors which is only possible if that something is a prime which we denote here by $p$. Note $p$ cannot be equal to $2$ nor $5$.

 Any other divisor of $10$ has to be $x$ but starting from $x = 5$ obviously doesn't work. We're done in this case and we get

 $\frac{2^4 \cdot 5(3 + 7 + 11 + 13 + 17 + 19 + 23)}{20} = 4(93) = 372$ in this case.

 Case 3: $k$ contains $5$ but not $2$ in its prime factorization

 This case is simple as we get $2^2 \cdot 5^{1 + a} \cdot \text{something}$ assuming $k$ has $a$ divisors of $5$. But the number of total divisors for this number is divisible by $3$ thanks to the $2^2$ term and $20$ is not divisible by $3$ so no cases here.

 Case 4: $k$ contains both $2$ and $5$ in its prime factorization

 Then the number is $2^{2 + a} \cdot 5^{1 + b} \cdot \text{something}$ assuming $k$ has $a$ factors of $2$ and $b$ factors of $5$. Assume that something has $x$ factors again. Then we get

 $(3 + a)(2 + b)x = 20$

 We will do casework noting $3 + a$ must be a divisor of $20$ that is at least $3$ since $a$ is nonnegative. Then we get

 $3 + a = 4, 5, 10, 20 \implies a = 1, 2, 7, 17$. We note $a = 17$ is clearly too big as the final number will be a multiple of $2^{19}$ which is clearly way too big. 

 $a = 7$ implies $(2 + b)x = 2 \implies b = 0, x = 1$ as the only solution that works if $x, b$ are nonnegative. But this doesn't work since this case assumes $k$ has both $2$ and $5$ in its prime factorization and if $b = 0$ then $k$ has no factors of $5$ a contradiction.

 $a = 2$ implies $(2 + b)x = 4 \implies x = 1, b = 2$ or $x = 2, b = 0$. As we know, $x = 2, b = 0$ can't work since $b = 0$ doesn't work. So $a = 2$ and $b = 2$ and $x = 1$ implies the number has to be $2^4 \cdot 5^3$. 

 $a = 1$ implies $(2 + b)x = 5 \implies x = 1, b = 3$. This gives $2^3 \cdot 5^4 = 5000$ which is too big.

 So for case 4, we only get one case and that is $2^4 \cdot 5^3 = 2000$.

 Dividing this by $20$ gives $100$.

 Finally, the answer is $372 + 100 = \boxed{472}$



	\item (\textit{AIME II 2019}) There is a unique angle $\theta$ between $0^\circ$ and $90^\circ$ such that for nonnegative integers $n,$ the value of $\tan(2^n\theta)$ is positive when $n$ is a multiple of $3$, and negative otherwise. The degree measure of $\theta$ is $\frac{p}{q}$, where $p$ and $q$ are relatively prime positive integers. Find $p+q$.


 
 \textbf{Answer: }547



\textbf{Solution 1}

Note that if $\tan \theta$ is positive, then $\theta$ is in the first or third quadrant, so $0^{\circ} < \theta < 90^{\circ} \pmod{180^{\circ}}$. 

 Furthermore, the only way $\tan{\left(2^{n}\theta\right)}$ can be positive for all $n$ that are multiples of $3$ is when:
 \[2^0\theta \equiv 2^3\theta \equiv 2^6\theta \equiv ... \pmod{180^{\circ}}.\](This is because if it isn't the same value, the terminal angle will gradually shift from the first quadrant into different quadrants, making the condition for positive tan untrue. This must also be true in order for $\theta$ to be unique.) 

 This is the case if $2^3\theta \equiv 2^0\theta \pmod{180^{\circ}}$, so $7\theta \equiv 0^{\circ} \pmod{180^{\circ}}$. Therefore, recalling that $0^{\circ}<\theta<90^{\circ},$ the possible $\theta$ are:
 \[\frac{180}{7}^{\circ}, \frac{360}{7}^{\circ}, \frac{540}{7}^{\circ}.\]
$\frac{180}{7}^{\circ}$ does not work since $\tan\left(2 \cdot \frac{180}{7}^{\circ}\right)$ is positive.

 $\frac{360}{7}^{\circ}$ does not work because $\tan\left(4 \cdot \frac{360}{7}^{\circ}\right)$ is positive.

 Thus, $\theta = \frac{540}{7}^{\circ}$, and a quick check verifies that it does work. Our desired answer is $540 + 7 = \boxed{547}$.

 



\textbf{Solution 2}

As in the previous solution, we note that $\tan \theta$ is positive when $\theta$ is in the first or third quadrant. In order for $\tan\left(2^n\theta\right)$ to be positive for all $n$ divisible by $3$, we must have $\theta$, $2^3\theta$, $2^6\theta$, etc to lie in the first or third quadrants. We already know that $\theta\in(0,90)$. We can keep track of the range of $2^n\theta$ for each $n$ by considering the portion in the desired quadrants, which gives 
 \[n=1 \implies (90,180)\]
 \[n=2\implies (270,360)\]
 \[n=3 \implies (180,270)\]
 \[n=4 \implies (90,180)\]
 \[n=5\implies(270,360)\]
 \[n=6 \implies (180,270)\]
 \[\cdots\]at which point we realize a pattern emerging. Specifically, the intervals repeat every $3$ after $n=1$. We can use these repeating intervals to determine the desired value of $\theta$ since the upper and lower bounds will converge to such a value (since it is unique, as indicated in the problem). Let's keep track of the lower bound.

 Initially, the lower bound is $0$ (at $n=0$), then increases to $\frac{90}{2}=45$ at $n=1$. This then becomes $45+\frac{45}{2}$ at $n=2$, $45+\frac{45}{2}$ at $n=3$, $45+\frac{45}{2}+\frac{45}{2^3}$ at $n=4$, $45+\frac{45}{2}+\frac{45}{2^3}+\frac{45}{2^4}$ at $n=5$, etc. Due to the observed pattern of the intervals, the lower bound follows a partial geometric series. Hence, as $n$ approaches infinity, the lower bound converges to 
 \[\sum_{k=0}^{\infty}\left(45+\frac{45}{2}\right)\cdot \left(\frac{1}{8}\right)^k=\frac{45+\frac{45}{2}}{1-\frac{1}{8}}=\frac{\frac{135}{2}}{\frac{7}{8}}=\frac{540}{7}\implies p+q=540+7=\boxed{547}\]-ktong

 \textbf{Extra note:} If you are still unsure, you can check the upper bound to see if the top converges. This value comes out to be: 

 $90-\frac{90}{2^3}-\frac{90}{2^6}-\cdots=90\left(1-\sum_{k=1}^{\infty}\frac{1}{2^{3k}}\right)=90\left(1-\frac{1}{7}\right)=\frac{540}{7}$.

 Thus, the limit does in fact converge to the same value on both sides. ~eevee9406

 \textbf{Extra extra note:}The top bound is arguably easier to prove because there are less terms added -dragoon

 



\textbf{Solution 3}

Since $\tan\left(\theta\right) > 0$, $0 < \theta < 90$.  Since $\tan\left(2\theta\right) < 0$, $\theta$ has to be in the second half of the interval (0, 90) ie (45, 90).  Since $\tan\left(4\theta\right) < 0$, $\theta$ has to be in the second half of that interval ie (67.5, 90).  And since $\tan\left(8\theta\right) > 0$, $\theta$ has to be in the first half of (67.5, 90).  Inductively, the pattern repeats: $\theta$ is in the second half of the second half of the first half of the second half of the second half... of the interval (0, 90).  Consider the binary representation of numbers in the interval (0, 1).  Numbers in the first half of the interval start with 0.0... and numbers in the second half start with 0.1... .  Similarly, numbers in the second half of the second half start with 0.11... etc.  So if we want a number in the first half of the second half of the second half... of the interval, we want its binary representation to be $0.11011011011..._2 = \frac{6}{7}_{10}$.  So we want the number which is 6/7 of the way through the interval (0, 90) so $\theta = \frac{6}{7}\cdot 90 = \frac{540}{7}$ and $p+q = 540 + 7 = \boxed{547}$




\textbf{Solution 4}

With some simple arithmetic and guess and check, we can set the lower bound and upper bounds for the "first round of $3$ powers of two", which are $\frac{540}{8} = 67.5$ and $\frac{630}{8} = 78.75$. Going on to the "second round of $3$ powers of two, we set the new lower and upper bounds as $\frac{360 \times 12.5}{64} = 70.3125$ and $\frac{360 \times 13.75}{64} = 77.34375$ using some guess and check and bashing. Now, it is obvious that the bounds for the "zeroth round of $3$ powers of two" are $0$ and $90$, and notice that $90 - 78.75 = 11.25$ and $78.75 - 77.34375 = 1.40625$ and $\frac{11.25}{1.40625} = 8$.This is obviously a geometric series, so setting $11.25$ as $u$, we obtain $90 - (u + \frac{u}{8} + \frac{u}{64} + ...)$ = $90 - \frac{u}{1-\frac{1}{8}}$ = $\frac{u}{\frac{7}{8}}$ = $\frac{45}{4} \times \frac{8}{7}$ which simplifies to $\frac{90}{7}$. We can now finally subtract $\frac{90}{7}$ from $\frac{630}{7}$ and then we get $\frac{540}{7}$ as the unique angle, so $\boxed{547}$ is our answer.-fidgetboss\_4000

 



\textbf{Solution 5}

Since $7$ is the only number n such that f(x) = $2^{\lfloor x\rfloor}$ $\text{(mod 7)}$  has a period of 3, we find that $\theta$ is a multiple of ${\frac{180}{7}}^\circ$. Note that the tangents of ${\frac{180}{7}}^\circ$, ${\frac{360}{7}}^\circ$, ${\frac{540}{7}}^\circ$ are positive while those of ${\frac{720}{7}}^\circ$, ${\frac{900}{7}}^\circ$, and ${\frac{1080}{7}}^\circ$ are negative. With a bit of trial and error, we find $\theta$ is ${\frac{540}{7}}^\circ$ and $540+7$ is $\boxed{547}$.

 


	\item (\textit{AIME II 2019}) Triangle $ABC$ has side lengths $AB=7,BC=8,$ and $CA=9.$ Circle $\omega_1$ passes through $B$ and is tangent to line $AC$ at $A.$ Circle $\omega_2$ passes through $C$ and is tangent to line $AB$ at $A.$ Let $K$ be the intersection of circles $\omega_1$ and $\omega_2$ not equal to $A.$ Then $AK=\frac{m}{n},$ where $m$ and $n$ are relatively prime positive integers. Find $m+n.$

 \textbf{Answer: }011



\textbf{Solution 1}


                    \begin{figure}[H]
                        \centering
                        \includegraphics[width=0.5\textwidth]{Solution/AIME_II_2019_11_Solution_Figure_1.png}
                    \end{figure}
                    Let $\tau$ be the $\sqrt{bc}$-inversion (also known as forced inversion overlay). By definition, $\tau$ sends $B\leftrightarrow C$. Using the tangency conditions, $\omega_1'$ is the line passing through $C$ parallel to $\overline{AB}$. Similarly, $\omega_2'$ is the line through $B$ parallel to $\overline{AC}$. It follows that $ABK'C$ is a parallelogram. 

 If $M$ is the midpoint of $BC$, then by Appolonius Theorem, $AM = 7$, so $AK'=14$. Thus, 
 \[AK = \frac{63}{AK'} = \frac{9}{2} \implies \boxed{011}\]

 -sillybone

 



\textbf{Solution 2}


                    \begin{figure}[H]
                        \centering
                        \includegraphics[width=0.5\textwidth]{Solution/AIME_II_2019_11_Solution_Figure_2.png}
                    \end{figure}
                    -Diagram by Brendanb4321

 
Note that from the tangency condition that the supplement of $\angle CAB$ with respects to lines $AB$ and $AC$ are equal to $\angle AKB$ and $\angle AKC$, respectively, so from tangent-chord, 
 \[\angle AKC=\angle AKB=180^{\circ}-\angle BAC\] Also note that $\angle ABK=\angle KAC$$^{(*)}$, so $\triangle AKB\sim \triangle CKA$. Using similarity ratios, we can easily find 
 \[AK^2=BK*KC\] However, since $AB=7$ and $CA=9$, we can use similarity ratios to get 
 \[BK=\frac{7}{9}AK, CK=\frac{9}{7}AK\]

            \begin{itemize}
        \item Now we use Law of Cosines on $\triangle AKB$: From reverse Law of Cosines, $\cos{\angle BAC}=\frac{11}{21}\implies \cos{(180^{\circ}-\angle BAC)}=\angle AKB=-\frac{11}{21}$
            \end{itemize}
            Giving us 
 \[AK^2+\frac{49}{81}AK^2+\frac{22}{27}AK^2=49\] 
 \[\implies \frac{196}{81}AK^2=49\] 
 \[AK=\frac{9}{2}\] so our answer is $9+2=\boxed{011}$.

 $^{(*)}$ Let $O$ be the center of $\omega_1$. Then $\angle KAC = 90 - \angle OAK = 90 - \frac{1}{2}(180 - \angle AOK) = \frac{\angle AOK}{2} = \angle ABK$. Thus, $\angle ABK = \angle KAC$

 -franchester; $^{(*)}$ by firebolt360

 \textbf{Supplement} 
            \begin{itemize}
        \item In order to get to the Law of Cosines first, we first apply the LOC to $\triangle{ABC},$ obtaining $\angle{BAC}.$
        \item We angle chase before applying the law of cosines to $\angle{AKB}.$
            \end{itemize}
            Note that $\angle{ABK}=\angle{KAC}$ and $\angle{KCA}=\angle{KAB}$ from tangent-chord. 

 Thus, $\angle{AKC}=\angle{AKB}=180^{\circ}-(\angle{ABK}+\angle{KAB}).$ 

 However from our tangent chord, note that:
 \[\angle{ABK}+\angle{KAB}=\angle{KAC}+\angle{KAB}=\angle{BAC}.\]Thus, $\angle{AKB}=180^\circ-\angle{BAC}.$

            \begin{itemize}
        \item As an alternative approach, note that the sum of the angles in quadrilateral $ABKC$ is $360^{\circ}$ and we can find $\angle{AKB}=\frac12$ of convex $\angle{BKC},$ which is just:
            \end{itemize}
            
 \[\frac12 \left(360^{\circ}-2(\angle{KAB}+\angle{KBA}\right) = 180^\circ - \angle{BAC}.\]




\textbf{Solution 3 (Inversion)}

Consider an inversion with center $A$ and radius $r=AK$. Then, we have $AB\cdot AB^*=AK^2$, or $AB^*=\frac{AK^2}{7}$. Similarly, $AC^*=\frac{AK^2}{9}$. Notice that $AB^*KC^*$ is a parallelogram, since $\omega_1$ and $\omega_2$ are tangent to $AC$ and $AB$, respectively. Thus, $AC^*=B^*K$. Now, we get that
 \[\cos(\angle AB^*K)=\cos(180-\angle BAC)=-\frac{11}{21}\]so by Law of Cosines on $\triangle AB^*K$ we have 
 \[(AK)^2=(AB^*)2+(B^*K)^2-2\cdot AB^*\cdot B^*K \cdot \cos(\angle AB^*K)\]
 \[\Rightarrow (AK)^2=\frac{AK^4}{49}+\frac{AK^4}{81}-2\cdot \frac{AK^2}{7}\frac{AK^2}{9}\frac{-11}{21}\]
 \[\Rightarrow 1=\frac{AK^2}{49}+\frac{AK^2}{81}+\frac{22AK^2}{63\cdot21}\]
 \[\Rightarrow AK=\frac{9}{2}\]Then, our answer is $9+2=\boxed{11}$. -brianzjk

 



\textbf{Solution 4 (Death By Trig Bash)}

Let the centers of the circles be $O_{1}$ and $O_{2}$ where the $O_{1}$ has the side length $7$ contained in the circle. Now let $\angle BAC =x.$ This implies 
 \[\angle O_{1}AB = \angle O_{1}BA = \angle O_{2}AC = \angle O_{2}CA = 90^{\circ}-x\] by the angle by by tangent. Then we also know that 
 \[\angle AO_{1}B = \angle AO_{2}C = 2x\] Now we first find $\cos x.$ We use law of cosines on $\bigtriangleup ABC$ to obtain 
 \[64 = 81 + 48 - 2 \cdot 9 \cdot 7 \cdot \cos{x}\] 
 \[\implies \cos{x} =\frac{11}{21}\] 
 \[\implies \sin{x} =\frac{8\sqrt{5}}{21}\] Then applying law of sines on $\bigtriangleup AO_{1}B$ we obtain 
 \[\frac{7}{\sin{2x}} =\frac{O_{1}B}{\sin{90^{\circ}-x}}\] 
 \[\implies\frac{7}{2\sin{x}\cos{x}} =\frac{O_{1}B}{\cos{x}}\] 
 \[\implies O_{1}B = O_{1}A=\frac{147}{16\sqrt{5}}\] Using similar logic we obtain $O_{2}A =\frac{189}{16\sqrt{5}}.$ 

 Now we know that $\angle O_{1}AO_{2}=180^{\circ}-x.$ Thus using law of cosines on $\bigtriangleup O_{1}AO_{2}$ yields 
 \[O_{1}O_{2} =\sqrt{\left(\frac{147}{16\sqrt{5}}\right)^2+\left(\frac{189}{16\sqrt{5}}\right)^2-2\:\cdot \left(\frac{147}{16\sqrt{5}}\right)\cdot \frac{189}{16\sqrt{5}}\cdot -\frac{11}{21}}\] While this does look daunting we can write the above expression as 
 \[\sqrt{\left(\frac{189+147}{16\sqrt{5}}\right)^2 - 2\cdot \left(\frac{147}{16\sqrt{5}}\right)\cdot \frac{189}{16\sqrt{5}}\cdot \frac{10}{21}} =\sqrt{\left(\frac{168}{8\sqrt{5}}\right)^2 - \left(\frac{7 \cdot 189 \cdot 5}{8 \sqrt{5} \cdot 8\sqrt{5}}\right)}\] Then factoring yields 
 \[\sqrt{\frac{21^2(8^2-15)}{(8\sqrt{5})^2}} =\frac{147}{8\sqrt{5}}\] The area 
 \[[O_{1}AO_{2}] =\frac{1}{2} \cdot\frac{147}{16\sqrt{5}} \cdot\frac{189}{16\sqrt{5}} \cdot \sin(180^{\circ}-x) =\frac{1}{2} \cdot\frac{147}{16\sqrt{5}} \cdot\frac{189}{16\sqrt{5}} \cdot\frac{8\sqrt{5}}{21}\] Now $AK$ is twice the length of the altitude of $\bigtriangleup O_{1}AO_{2}$ so we let the altitude be $h$ and we have 
 \[\frac{1}{2} \cdot h \cdot\frac{147}{8\sqrt{5}} =\frac{1}{2} \cdot\frac{147}{16\sqrt{5}} \cdot\frac{189}{16\sqrt{5}} \cdot\frac{8\sqrt{5}}{21}\] 
 \[\implies h =\frac{9}{4}\] Thus our desired length is $\frac{9}{2} \implies m+n = \boxed{11}.$

 -minor edits by faliure167

 


	\item (\textit{AIME II 2019}) For $n\ge1$ call a finite sequence $(a_1,a_2,\ldots,a_n)$ of positive integers \textit{progressive} if $a_i<a_{i+1}$ and $a_i$ divides $a_{i+1}$ for $1\le i\le n-1$. Find the number of progressive sequences such that the sum of the terms in the sequence is equal to $360.$

 \textbf{Answer: }047



\textbf{Solution}

If the first term is $x$, then dividing through by $x$, we see that we can find the number of progressive sequences whose sum is $\frac{360}{x} - 1$, and whose first term is not 1. If $a(k)$ denotes the number of progressive sequences whose sum is $k$ and whose first term is not 1, then we can express the answer $N$ as follows:

 
 \begin{align*}N &= a(359) + a(179) + a(119) + a(89) + a(71) + a(59) + a(44) + a(39) \\ &+ a(35) + a(29) + a(23) + a(19) + a(17) + a(14) + a(11) + a(9) \\ &+ a(8) + a(7) + a(5) + a(4) + a(3) + a(2) + a(1) + 1 \end{align*}


 And in general, we have $a(k) = 1+\sum_{d|k\text{ and } d\ne 1,k}a(d-1)$.

 The $+1$ at the end accounts for the sequence whose only term is 360. Fortunately, many of these numbers are prime; we have $a(p) = 1$ for primes $p$ as the only such sequence is "$p$" itself. Also, $a(1) = 0$. So we have

 
 \[N = 15 + a(119) + a(44) + a(39) + a(35) + a(14) + a(9) + a(8) + a(4)\]

 For small $k$, $a(k)$ is easy to compute: $a(4) = 1$, $a(8) = 2$, $a(9) = 2$. For intermediate $k$ (e.g. $k=21$ below), $a(k)$ can be computed recursively using previously-computed values of $a(\cdot)$, similar to dynamic programming. Then we have
 \begin{align*} a(14) &= 1+a(6) = 1+2 = 3\\ a(35) &= 1+a(6)+a(4) = 1 + 2 + 1 = 4 \\ a(39) &= 2 + a(12) = 2+4 = 6 \\ a(44) &= 2 + a(21) + a(10) = 2+4+2=8 \\ a(119) &= 1 + a(16) + a(6) = 1 + 3 + 2 = 6 \end{align*}
Thus the answer is $N = 15+6+8+6+4+3+2+2+1 = \boxed{047}$.

 -scrabbler94

 


	\item (\textit{AIME II 2019}) Regular octagon $A_1A_2A_3A_4A_5A_6A_7A_8$ is inscribed in a circle of area $1.$ Point $P$ lies inside the circle so that the region bounded by $\overline{PA_1},\overline{PA_2},$ and the minor arc $\widehat{A_1A_2}$ of the circle has area $\frac{1}{7},$ while the region bounded by $\overline{PA_3},\overline{PA_4},$ and the minor arc $\widehat{A_3A_4}$ of the circle has area $\frac{1}{9}.$ There is a positive integer $n$ such that the area of the region bounded by $\overline{PA_6},\overline{PA_7},$ and the minor arc $\widehat{A_6A_7}$ of the circle is equal to $\frac{1}{8}-\frac{\sqrt2}{n}.$ Find $n.$

 \textbf{Answer: }504



\textbf{Solution 1}

The actual size of the diagram doesn't matter. To make calculation easier, we discard the original area of the circle, \(1\), and assume the side length of the octagon is \(2\). 

 Let \(r\) denote the radius of the circle, \(O\) be the center of the circle. Then:
 \[r^2= 1^2 + \left(\sqrt{2}+1\right)^2= 4+2\sqrt{2}.\]

 Now, we need to find the "D" shape, the small area enclosed by one side of the octagon and \(\frac{1}{8}\) of the circumference of the circle: 
 \[D= \frac{1}{8} \pi r^2 - [A_1 A_2 O]=\frac{1}{8} \pi \left(4+2\sqrt{2}\right)- \left(\sqrt{2}+1\right)\]

 Let \(PU\) be the height of \(\triangle A_1 A_2 P\), \(PV\) be the height of \(\triangle A_3 A_4 P\), \(PW\) be the height of \(\triangle A_6 A_7 P\). From the \(\frac{1}{7}\) and \(\frac{1}{9}\) condition we have 
 \[\triangle P A_1 A_2= \frac{\pi r^2}{7} - D= \frac{1}{7} \pi \left(4+2\sqrt{2}\right)-\left(\frac{1}{8} \pi \left(4+2\sqrt{2}\right)- \left(\sqrt{2}+1\right)\right)\]
 \[\triangle P A_3 A_4= \frac{\pi r^2}{9} - D= \frac{1}{9} \pi \left(4+2\sqrt{2}\right)-\left(\frac{1}{8} \pi \left(4+2\sqrt{2}\right)- \left(\sqrt{2}+1\right)\right)\]which gives \(PU= \left(\frac{1}{7}-\frac{1}{8}\right) \pi \left(4+ 2\sqrt{2}\right) + \sqrt{2}+1\) and \(PV= \left(\frac{1}{9}-\frac{1}{8}\right) \pi \left(4+ 2\sqrt{2}\right) + \sqrt{2}+1\). 

 Now, let \(A_1 A_2\) intersects \(A_3 A_4\) at \(X\), \(A_1 A_2\) intersects \(A_6 A_7\) at \(Y\),\(A_6 A_7\) intersects \(A_3 A_4\) at \(Z\). 

 Clearly, \(\triangle XYZ\) is an isosceles right triangle, with right angle at \(X\) and the height with regard to which shall be \(3+2\sqrt2\). Now \(\frac{PU}{\sqrt{2}} + \frac{PV}{\sqrt{2}} + PW = 3+2\sqrt2\) which gives:

 \begin{align*}PW&= 3+2\sqrt2-\frac{PU}{\sqrt{2}} - \frac{PV}{\sqrt{2}} \\&=3+2\sqrt{2}-\frac{1}{\sqrt{2}}\left(\left(\frac{1}{7}-\frac{1}{8}\right) \pi \left(4+ 2\sqrt{2}\right) + \sqrt{2}+1+\left(\frac{1}{9}-\frac{1}{8}\right) \pi \left(4+ 2\sqrt{2}\right) + \sqrt{2}+1\right)\\&=1+\sqrt{2}- \frac{1}{\sqrt{2}}\left(\frac{1}{7}+\frac{1}{9}-\frac{1}{4}\right)\pi\left(4+2\sqrt{2}\right)\end{align*}


 Now, we have the area for \(D\) and the area for \(\triangle P A_6 A_7\), so we add them together:

 \begin{align*}\text{Target Area} &= \frac{1}{8} \pi \left(4+2\sqrt{2}\right)- \left(\sqrt{2}+1\right) + \left(1+\sqrt{2}\right)- \frac{1}{\sqrt{2}}\left(\frac{1}{7}+\frac{1}{9}-\frac{1}{4}\right)\pi\left(4+2\sqrt{2}\right)\\&=\left(\frac{1}{8} - \frac{1}{\sqrt{2}}\left(\frac{1}{7}+\frac{1}{9}-\frac{1}{4}\right)\right)\text{Total Area}\end{align*}


 The answer should therefore be \(\frac{1}{8}- \frac{\sqrt{2}}{2}\left(\frac{16}{63}-\frac{16}{64}\right)=\frac{1}{8}- \frac{\sqrt{2}}{504}\). The answer is \(\boxed{504}\).

 \textbf{SpecialBeing2017}

 



\textbf{Solution 2}

Instead of considering the actual values of the areas, consider only the changes in the areas that result from moving point $P$ from the center of the circle. We will proceed by coordinates. Set the origin at the center of the circle and refer to the following diagram, where the octagon is oriented so as $A_1A_2$ is horizontal (and therefore $A_3A_4$ is vertical). Note that the area bounded by $\overline{A_iA_j}$ and the arc $\widehat{A_iA_j}$ is fixed, so we only need to consider the relevant triangles.

 
                    \begin{figure}[H]
                        \centering
                        \includegraphics[width=0.5\textwidth]{Solution/AIME_II_2019_13_Solution_Figure_1.png}
                    \end{figure}
                    

 Define one arbitrary unit as the distance that you need to move $P$ from $A_1A_2$ to change the area of $\triangle PA_1A_2$ by $1$. We can see that $P$ was moved down by $\frac{1}{7}-\frac{1}{8}=\frac{1}{56}$ units to make the area defined by $P$, $A_1$, and $A_2$ $\frac{1}{7}$. Similarly, $P$ was moved right by $\frac{1}{8}-\frac{1}{9}=\frac{1}{72}$ to make the area defined by $P$, $A_3$, and $A_4$ $\frac{1}{9}$. This means that $P$ has coordinates $(\frac{1}{72},-\frac{1}{56})$.

 Now, we need to consider how this displacement in $P$ affected the area defined by $P$, $A_6$, and $A_7$. This is equivalent to finding the shortest distance between $P$ and the blue line in the diagram (as $K=\frac{1}{2}bh$ and the blue line represents $h$ while $b$ is fixed). Using an isosceles right triangle, one can find the that shortest distance between $P$ and this line is $\frac{\sqrt{2}}{2}(\frac{1}{56}-\frac{1}{72})=\frac{\sqrt{2}}{504}$.

 Remembering the definition of our unit, this yields a final area of
 \[\frac{1}{8}-\frac{\sqrt{2}}{\boxed{504}}.\]

 -Archeon

 



\textbf{Solution 3 (Shoelace)}

Let $r=1/\sqrt{\pi}$ be the radius of the circle, and plot the octagon on the Cartesian plane such that $A_i = (r\cos(\pi i/4),r\sin(\pi i/4))$ for each $i$. Let $P$ have coordinates $(x,y)$. Note that the area bounded by the circle and any segment of the form $A_iA_{i+1}$ is 
 \[S=\frac{1}{8}-\frac{1}{2}r^2\sin\left(\frac{\pi}{4}\right) = \frac{1}{8}-\frac{r^2}{\sqrt{8}}.\]According to the Shoelace formula, we have \begin{align*}[A_1A_2P] &= \frac{1}{2}\left(\frac{r^2}{\sqrt{2}}+\frac{rx}{\sqrt{2}} - rx - \frac{ry}{\sqrt{2}}\right) = \frac{r^2+rx(1-\sqrt{2})-ry}{\sqrt{8}},\\[A_3A_4P] &= \frac{1}{2}\left(-ry+\frac{rx}{\sqrt{2}} + \frac{r^2}{\sqrt{2}} + \frac{ry}{\sqrt{2}}\right) = \frac{r^2+rx+ry(1-\sqrt{2})}{\sqrt{8}}.\end{align*}
Note that $[A_1A_2P] = \frac{1}{7}-S$ and $[A_3A_4P] = \frac{1}{9}-S$. Let $X=rx$ and $Y=ry$. Using the work above, we obtain the system\begin{align*}X(1-\sqrt{2}) - Y &= \frac{\sqrt{8}}{56}\\X+Y(1-\sqrt{2}) &= -\frac{\sqrt{8}}{72},\end{align*}


 which yields the solution 
 \[X = -\frac{1}{56}-\frac{1+\sqrt{2}}{72},\quad Y=-\frac{1+\sqrt{2}}{56}+\frac{1}{72}.\]Then by the Shoelace formula, the area bounded by $PA_6,PA_7$ and arc $A_6A_7$ is\begin{align*}[A_6A_7P]+S &= \frac{r^2+rx(1-\sqrt{2})+ry}{8} + \frac{1}{8} - \frac{r^2}{\sqrt{8}}\\&= \frac{1}{8} + \frac{1-\sqrt{2}}{\sqrt{8}}\left(-\frac{1}{56}-\frac{1+\sqrt{2}}{72}\right)  + \frac{1}{\sqrt{8}}\left(-\frac{1+\sqrt{2}}{56}+\frac{1}{72}\right)\\&= \frac{1}{8} +\frac{1}{72\sqrt{2}} - \frac{1}{56\sqrt{2}}\\&= \frac{1}{8} - \frac{\sqrt{2}}{504}.\end{align*}
The requested integer is $n=\boxed{504}.$~TThB0501

 


	\item (\textit{AIME II 2019}) Find the sum of all positive integers $n$ such that, given an unlimited supply of stamps of denominations $5,n,$ and $n+1$ cents, $91$ cents is the greatest postage that cannot be formed.


 
 \textbf{Answer: }071



\textbf{Solution 1}

We will do casework on $n \pmod{5}$ and see which values work. 

 Case 1: $n \equiv 1 \pmod{5}$
$n$ would have to be greater than $91$ because if it wasn't then we could just add denominations of $5$ until we got to $91.$ However, if $n$ is more than $91,$ then all the numbers from $91$ to $n$ which are not multiples of $5$ are also unrepresentable, which violates the problem's condition of $91$ being the largest unrepresentable number. So there are no values for this case. 

 
Case 2: $n \equiv 2 \pmod{5}$
$n + 1 \equiv 3 \pmod{5}$ is the smallest formable number 3 mod 5, so $n + 1 - 5 = n - 4$ is the largest unformable number 3 mod 5. $2n + 2 \equiv 1 \pmod{5}$ is the smallest formable number 1 mod 5, so $2n + 2 - 5 = 2n - 3$ is the largest unformable number 1 mod 5. Finally, $2n \equiv 4 \pmod{5}$ is the smallest formable number 4 mod 5, so $2n - 5$ is the largest unformable number 4 mod 5.

 Therefore, the largest unformable number overall is $2n - 3 \equiv 1 \pmod{5}.$ This satisfies the problem's condition!  $2n - 3 = 91 \implies n = 47.$ 

 
Case 3: $n \equiv 3 \pmod{5}$
$n + 1 \equiv 4 \pmod{5},$ which is the smallest representable number 4 mod 5, so the largest unrepresentable number 4 mod 5 is $n + 1 - 5 = n -4.$ 

 Similarly, $2n \equiv 1 \pmod{5}$ is the smallest formable number 1 mod 5 and $2n + 1 \equiv 2 \pmod{5}$ (Yes ~Towerrocks, it was 2n + 1 my bad!) is the smallest formable number 2 mod 5, so the largest number 2 mod 5 that is unrepresentable is $2n + 1 - 5 = 2n - 4,$ but the largest number 1 mod 5 that is unrepresentable is $2n -5.$ $2n - 4 > 2n - 5,$ so there would be a larger unrepresentable number than $91$ if we were to set $2n - 5$ equal to it, which violates the problem's condition. And we can't set $2n - 4$ equal to $91$ because $91 \not\equiv 2 \pmod{5},$ so there are no valid values for this case.

 


 Case 4: $n \equiv 4 \pmod{5}$
$n + 1$ would be a multiple of $5,$ so we would be left with denominations of $5$ and $n.$ Then it would just be a straightforward application of the Chicken Mcnugget Theorem: $5n - n - 5 = 91 \implies n = 24,$ which is indeed $4 \pmod{5}.$ So this value works.

 
Case 5: $n$ is a multiple of 5 

 In this case we’re left with denominations of $5$ and $n + 1,$ so by the Chicken Mcnugget Theorem we get $5(n + 1) - n - 6 = 91 \implies n = 23,$ which is not a multiple of $5.$ So there are no valid values here.

 
The two valid values we found were $47$ and $24,$ so the answer is $47 + 24 = \boxed{71}.$




\textbf{Solution 2}

Notice that if we let $n + 1$ be a multiple of $5,$ then by the Chicken McNugget Theorem we have $5n - (n + 5) = 91 \implies n = 24.$ We find this is the smallest possible value of $n$.

 For a value of $n$ to work, we should not be able to form the value $91,$ but we should be able to form the values from $92$ to $96$. After $96,$ we can form any value by adding additional $5$ cent stamps ($92 + 5x$, $93 + 5x$, $94 + 5x$, $95 + 5x$, and $96 + 6x$, where $x$ is a positive integer, together make all positive integers after $96$). We also need to be able to form $96$ using only denominations of $n$ and $n + 1,$ because if we also used denominations of $5,$ then we could just remove a $5$ cent stamp to get back to $91.$

 Now we can figure out the working $(n, n+1)$ pairs that can be used to obtain $96$. We also need to make sure we can get the numbers $92 - 95$ but not $91$ with denominations of $5$. We can use no more than a total of $\frac{96}{24}=4$ stamps. Let $an + b(n + 1) = 96,$ where $a$ is the number of $n$ stamps used and $b$ is the number of $n + 1$ stamps used. The possible $(a,b)$ pairs are:
 \[(a,b) \rightarrow (4,0), (3,1), (2,2), (1,3), (0,4), (3,0), (2,1), (1,2), (0,3), (2,0), (1,1), (0,2), (1,0), (0,1).\]Solving for $n$ in each $(a,b)$ pair we find that the potential solutions are:
 \[(n, n + 1) \rightarrow (23,24), (24, 25), (31, 32), (32, 33), (47, 48), (48, 49), (95, 96), (96, 97).\]

 The first pair doesn't work because $n \geq 24,$ and the last two don't work either because they are too large to form the values $92$ through $95$. Finally, $(31,32)$ and $(32, 33)$ don't work because they can also form $91$ with denominations of $5$ (for example, $32 \cdot 0 + 33\cdot 2 + 5 \cdot 5 = 91$). We are left with $(24, 25), (47, 48), (48,49).$ Now we need to make sure that $92 - 95$ can be formed. We know from the Chicken McNugget Theorem that $91$ is the largest unformable number with denominations $24, 25$ and $5,$ so this means that all numbers after $91$ should be formable, and we can bash to confirm that all numbers from $92-95$ can be formed with denominations of $47, 48,$ and $5.$ However, we find that $92$ cannot be formed with denominations of $48, 49,$ and $5,$ so $n = 48$ is eliminated. Then $n \in \{24, 47\}$. The requested sum is $24 + 47 = \boxed{071}$.

 \textbf{Note on finding and testing potential pairs} In order to find potential $(n,n+1)$ pairs, we simply test all combinations of $n$ and $n+1$ that sum to less than $4n$ (so that $n\ge24$) to see if they produce an integer value of $n$ when their sum is set to $96$. Note that, since $96$ is divisible by $1$, $2$, $3$, and $4$, we must either use only $n$ or only $n+1$, as otherwise, the sum is guaranteed to not be divisible by one of the numbers $2$, $3$, and $4$.

 $\begin{array}{c|c|c} \text{Combination} & \text{Sum} & n\text{-value} \\ \hline n,n,n,n & 4n & 24 \\ n+1,n+1,n+1 & 3n+3 & 31 \\ n,n,n & 3n & 32 \\ n+1,n+1 & 2n+2 & 47 \\ n,n & 2n & 48 \\ n+1 & n+1 & 95 \\ n & n & 96 \\ \end{array}$

 To test whether a pair works, we simply check that, using the number $5$ and the two numbers in the pair, it is impossible to form a sum of $91$, and it is possible to form sums of $92$, $93$, and $94$. ($95$ can always be formed using only $5$s, and the pair is already able to form $96$ because that was how it was found.) We simply need to reach the residues $2$, $3$, and $4$$\pmod{5}$ using only $n$ and $n+1$ without going over the number we are trying to form, while being unable to do so with the residue $1$. As stated in the above solution, the last two pairs are clearly too large to work.

 $\begin{array}{c|c|c|c|c} \text{Pair} & \text{Not }91 & 92 & 93 & 94 \\ \hline 24,25 & \checkmark & \checkmark & \checkmark & \checkmark \\ 31,32 & \times & \checkmark & \checkmark & \checkmark \\ 32,33 & \times & \checkmark & \checkmark & \checkmark \\ 47,48 & \checkmark & \checkmark & \checkmark & \checkmark \\ 48,49 & \checkmark & \times & \checkmark & \checkmark \\ \end{array}$

 (Note that if a pair is unable to fulfill a single requirement, there is no need to check the rest.)

 



\textbf{Solution 3}

Notice that once we hit all residues $\bmod 5$, we'd be able to get any number greater (since we can continually add $5$ to each residue). Furthermore, $n\not\equiv 0,1\pmod{5}$ since otherwise $91$ is obtainable (by repeatedly adding $5$ to either $n$ or $n+1$) Since the given numbers are $5$, $n$, and $n+1$, we consider two cases: when $n\equiv 4\pmod{5}$ and when $n$ is not that. 

 When $n\equiv 4 \pmod{5}$, we can only hit all residues $\bmod 5$ once we get to $4n$ (since $n$ and $n+1$ only contribute $1$ more residue $\bmod 5$). Looking at multiples of $4$ greater than $91$ with $n\equiv 4\pmod{5}$, we get $n=24$. It's easy to check that this works. Furthermore, any $n$ greater than this does not work since $91$ isn't the largest unobtainable value (can be verified using Chicken McNugget Theorem). 

 Now, if $n\equiv 2,3\pmod{5}$, then we'd need to go up to $2(n+1)=2n+2$ until we can hit all residues $\bmod 5$ since $n$ and $n+1$ create $2$ distinct residues $\bmod{5}$. Checking for such $n$ gives $n=47$ and $n=48$. It's easy to check that $n=47$ works, but $n=48$ does not (since $92$ is unobtainable). Furthermore, any $n$ greater than this does not work since $91$ isn't the largest unobtainable value in those cases (can be verified using Chicken McNugget Theorem). (Also note that in the $3 \pmod{5}$ case, the residue $2 \pmod{5}$ has will not be produced until $3(n+1)$ while the $1\pmod5$ case has already been produced, so the highest possible value that cannot be produced would not be a number equivalent to $1 \pmod5$)

 Since we've checked all residues $\bmod 5$, we can be sure that these are all the possible values of $n$. Hence, the answer is $24+47=\boxed{071}$. - ktong

 



\textbf{Solution 4}

Obviously $n\le 90$. We see that the problem's condition is equivalent to: 96 is the smallest number that can be formed which is 1 mod 5, and 92, 93, 94 can be formed (95 can always be formed). Now divide this up into cases. If $n\equiv 0\pmod{5}$, then 91 can be formed by using $n+1$ and some 5's, so there are no solutions for this case. If $n\equiv 1\pmod{5}$, then 91 can be formed by using $n$ and some 5's, so there are no solutions for this case either.

 For $n\equiv 2\pmod{5}$, $2n+2$ is the smallest value that can be formed which is 1 mod 5, so $2n+2=96$ and $n=47$. We see that $92=45+47$, $93=48+45$, and $94=47+47$, so $n=47$ does work. If $n\equiv 3\pmod{5}$, then the smallest value that can be formed which is 1 mod 5 is $2n$, so $2n=96$ and $n=48$. We see that $94=49+45$ and $93=48+45$, but 92 cannot be formed, so there are no solutions for this case. If $n\equiv 4\pmod{5}$, then we can just ignore $n+1$ since it is a multiple of 5, meaning that the Chicken McNugget theorem is a both necessary and sufficient condition, and it states that $5n-n-5=91$ meaning $4n=96$ and $n=24$.Hence, the only two $n$ that work are $n=24$ and $n=47$, so our answer is $24+47=\boxed{071}$.-Stormersyle

 



\textbf{Solution 5 (standard)}

Consider a postage that gives $96$. We cannot use a $5$-stamp as otherwise simply removing it yields a postage that gives $91$. Additionally, there cannot be at least $5$ of $n$-stamps or $n + 1$-stamps, as else we can convert $5$ of the same valued stamp into a positive number of $5$-stamps, then remove one to get a postage of $91$.

 
From here consider integers $0 \le a, b, \le 4$ where $an + b(n + 1) = 96 \implies n = \dfrac{96 - b}{a + b}$. The only pairs $(a, b)$ that yield an integer value are $(a, b) = (0, 2), (0, 3), (0, 4), (1, 0), (2, 0), (3, 0), (4, 0), (4, 1)$ which generate the values $n = 47, 31, 23, 96, 48, 32, 24, 19$ respectively. It is easy to find counterexamples of postages that evaluate to $91$ besides $n = 24, 47$. 

 
Now for $n = 24$ clearly $91$ is unobtainable since we need a $4 \pmod {5}$ amount of $24$-stamps which exceeds a value of $96$. A similar result holds for $n = 47$ as any evaluation $\le 2 \cdot 47$ can only be $0, 2, 3 \pmod{5}$. In both cases it is easy to construct a postage for $92, 93, 94, 95, 96$, to which repeatedly adding $5$-stamps makes all postages worth $>91$ possible. The requesteed sum is $24 + 47 = \boxed{71}$.

 - blueprimes

 



\textbf{Solution 6}

We can use the formula $g = \left( \max_{j \in \{1, 2, \ldots, a_1-1\}} R_j \right) - a_1$, where $g$ is the Frobenius Number (in this case 91), $a_1$ is the smallest denomination (which is 5), $j \in \{1, 2, \ldots, a_1-1\}$ represents a remainder (residue class) when you divide an integer by $a_1,$ and $R_j$ is the \textbf{smallest} nonnegative integer such that $a \cdot n + b \cdot (n + 1) \equiv j \pmod{5}$ where $a$ and $b$ are also nonnegative integers. So we have $91 = \left( \max_{j \in \{1, 2, 3, 4\}} R_j \right) - 5 \implies 96 = \left( \max_{j \in \{1, 2, 3, 4\}} R_j \right).$ This means that the largest of the $R_j$ values must be equal to exactly $96.$ We will do casework on $n \pmod{5}.$

 
\textbf{Case 1:} $n \equiv 1 \pmod{5}.$
$n + 1 \equiv 2 \pmod{5}.$ Plug this in:

 $R_1$: $a \cdot 1 + b \cdot 2 \equiv 1 \pmod{5} \implies a = 1, b = 0.$ So $R_1 = 1 \cdot n + 0 \cdot (n + 1) = n.$
$R_2$: $a \cdot 1 + b \cdot 2 \equiv 2 \pmod{5} \implies a = 0, b = 1.$ So $R_2 = 0 \cdot n + 1 \cdot (n + 1) = n + 1.$
$R_3$: $a \cdot 1 + b \cdot 2 \equiv 3 \pmod{5} \implies a = 1, b = 1.$ So $R_3 = 1 \cdot n + 1 \cdot (n + 1) = 2n + 1.$
$R_4$: $a \cdot 1 + b \cdot 2 \equiv 4 \pmod{5} \implies a = 0, b = 2.$ So $R_4 = 2n + 2.$ This is the largest of the four $R_j$ values. Setting this equal to $96$ we get $n = 47.$ But this doesn't work because $47$ is not congruent to $1 \pmod{5}.$ 

 
\textbf{Case 2:} $n \equiv 2 \pmod{5}.$
$n + 1 \equiv 3 \pmod{5}.$ Plug this in:

 $R_1$: $a \cdot 2 + b \cdot 3 \equiv 1 \pmod{5} \implies a = 0, b = 2.$ So $R_1 = 0 \cdot n + 2 \cdot (n + 1) = 2n + 2.$
$R_2$: $a \cdot 2 + b \cdot 3 \equiv 2 \pmod{5} \implies a = 1, b = 0.$ So $R_2 = 1 \cdot n + 0 \cdot (n + 1) =  n.$
$R_3$: $a \cdot 2 + b \cdot 3 \equiv 3 \pmod{5} \implies a = 0, b = 1.$ So $R_3 = 0 \cdot n + 1 \cdot (n + 1) = n + 1.$
$R_4$: $a \cdot 2 + b \cdot 3 \equiv 4 \pmod{5} \implies a = 2, b = 0.$ So $R_4 = 2 \cdot n + 0 \cdot (n + 1) = 2n.$ 

 The largest of the four $R_j$ values is $R_1 = 2n + 2.$ Setting this equal to $96$ we get $2n + 2 = 96 \implies n = 47.$ Now this works! $47$ is indeed $2 \pmod{5}.$

 
\textbf{Case 3:} $n \equiv 3 \pmod{5}.$
$n + 1 \equiv 4 \pmod{5}.$ Plug this in:

 $R_1$: $a \cdot 3 + b \cdot 4 \equiv 1 \pmod{5} \implies a = 2, b = 0.$ So $R_1 = 2 \cdot n + 0 \cdot (n + 1) = 2n.$
$R_2$: $a \cdot 3 + b \cdot 4 \equiv 2 \pmod{5} \implies a = 1, b = 1.$ So $R_2 = 1 \cdot n + 1 \cdot (n + 1) = 2n + 1.$
$R_3$: $a \cdot 3 + b \cdot 4 \equiv 3 \pmod{5} \implies a = 1, b = 0.$ So $R_3 = 1 \cdot n + 0 \cdot (n + 1) = n.$
$R_4$: $a \cdot 3 + b \cdot 4 \equiv 4 \pmod{5} \implies a = 0, b = 1.$ So $R_4 = 0 \cdot n + 1 \cdot (n + 1) = n + 1.$ 

 The largest of the four $R_j$ values is $R_2 = 2n + 1.$ Setting this equal to $96$ we get $n = 47.5,$ which is not an integer. So this doesn't work.

 
\textbf{Case 4:} $n \equiv 4 \pmod{5}.$ 

 We notice that $n + 1$ will be a multiple of $5,$ so we're left with the denominations of $5$ and $n$. Since these are two relatively prime numbers, we can do a straightforward application of the Chicken Mcnugget Theorem to get $5n - 5 - n = 91 \implies n = 24.$ This works because $24 \equiv 4 \pmod{5}.$ 

 
\textbf{Case 5:} $n \equiv 0 \pmod{5}.$

 Same thing as case $4,$ only this time we are left with denominations of $5$ and $n + 1.$ Applying the Chicken Mcnugget Theorem gives $4n - 1 = 91 \implies n = 23,$ but this doesn't work because $23$ clearly isn't a multiple of $5$.  

 The two valid values we found were $47$ and $24,$ so the sum is $\boxed{71}.$



	\item (\textit{AIME II 2019}) In acute triangle $ABC,$ points $P$ and $Q$ are the feet of the perpendiculars from $C$ to $\overline{AB}$ and from $B$ to $\overline{AC}$, respectively. Line $PQ$ intersects the circumcircle of $\triangle ABC$ in two distinct points, $X$ and $Y$. Suppose $XP=10$, $PQ=25$, and $QY=15$. The value of $AB\cdot AC$ can be written in the form $m\sqrt n$ where $m$ and $n$ are positive integers, and $n$ is not divisible by the square of any prime. Find $m+n$.


 
 \textbf{Answer: }574



\textbf{Solution 1}

First we have $a\cos A=PQ=25$, and $(a\cos A)(c\cos C)=(a\cos C)(c\cos A)=AP\cdot PB=10(25+15)=400$ by PoP. Similarly, $(a\cos A)(b\cos B)=15(10+25)=525,$ and dividing these each by $a\cos A$ gives$b\cos B=21,c\cos C=16$. 

 It is known that the sides of the orthic triangle are $a\cos A,b\cos B,c\cos C$, and its angles are $\pi-2A$,$\pi-2B$, and $\pi-2C$. We thus have the three sides of the orthic triangle now.Letting $D$ be the foot of the altitude from $A$, we have, in $\triangle DPQ$, 
 \[\cos P,\cos Q=\frac{21^2+25^2-16^2}{2\cdot 21\cdot 25},\frac{16^2+25^2-21^2}{2\cdot 16\cdot 25}= \frac{27}{35}, \frac{11}{20}.\]
 \[\Rightarrow \cos B=\cos\left(\frac 12 (\pi-P)\right)=\sin\frac 12 P =\sqrt{\frac{4}{35}},\]similarly, we get  
 \[\cos C=\cos\left(\frac 12 (\pi-Q)\right)=\sin\frac 12 Q=\sqrt{\frac{9}{40}}.\]To finish, 
 \[bc= \frac{(b\cos B)(c\cos C)}{\cos B\cos C}=\frac{16\cdot 21}{(2/\sqrt{35})(3/\sqrt{40})}=560\sqrt{14}.\]The requested sum is $\boxed{574}$.- crazyeyemoody907

 Remark: The proof that $a \cos A = PQ$ can be found here: \href{https://artofproblemsolving.comhttp://www.irmo.ie/5.Orthic\_triangle.pdf}{http://www.irmo.ie/5.Orthic\_triangle.pdf}

 



\textbf{Solution 2}

Let $BC=a$, $AC=b$, and $AB=c$. Let $\cos\angle A=k$. Then $AP=bk$ and $AQ=ck$. 

 By Power of a Point theorem, 
 \begin{align}     AP\cdot BP=XP\cdot YP \quad &\Longrightarrow \quad b^2k^2-bck+400=0\\     AQ\cdot CQ=YQ\cdot XQ \quad &\Longrightarrow \quad c^2k^2-bck+525=0 \end{align}Thus $bck = (bk)^2+400=(ck)^2+525 = u$. Then $bk=\sqrt{u-400}$, $ck=\sqrt{u-525}$, and 
 \[k=\sqrt{\frac{(u-400)(u-525)}{u^2}}\]Use the Law of Cosines in $\triangle APQ$ to get $25^2=b^2k^2+c^2k^2-2bck^3 = 2bck-925-2bck^3$, which rearranges to 
 \[775=bck - k^2\cdot bck  = u-\frac{(u-400)(u-525)}{u}\]Upon simplification, this reduces to a linear equation in $u$, with solution $u=1400$. Then 
 \[AB\cdot AC = bc = \frac 1{k}\cdot bck = \frac{u^2}{\sqrt{(u-400)(u-525)}}=560 \sqrt{14}\]So the final answer is $560 + 14 = \boxed{574}$

 By SpecialBeing2017

 



\textbf{Solution 3}

Let $AP=p$, $PB=q$, $AQ=r$, and $QC=s$.  By Power of a Point, 
 \begin{align}  AP\cdot PB=XP\cdot YP \quad &\Longrightarrow \quad pq=400\\  AQ\cdot QC=YQ\cdot XQ \quad &\Longrightarrow \quad rs=525  \end{align}Points $P$ and $Q$ lie on the circle, $\omega$, with diameter $BC$, and pow$(A,\omega) = AP\cdot AB = AQ\cdot AC$, so 
 \[p(p+q)=r(r+s)\quad \Longrightarrow \quad p^2-r^2=125\] Use Law of Cosines in $\triangle APQ$ to get $25^2=p^2+r^2-2pr\cos A$; since $\cos A = \frac r{p+q}$, this simplifies as
 \[500 \ =\  2r^2-\frac{2pr^2}{p+q} \ =\  2r^2-\frac{2p^2r^2}{p^2+400} \ =\ \frac{800r^2}{r^2+525}\]We get $r=5\sqrt{35}$ and thus
 \[r=5\sqrt{35}, \quad p = \sqrt{r^2+125} = 10\sqrt{10}, \quad q = \frac{400}{p} =4\sqrt{10}, \quad s= \frac{525}{r} = 3\sqrt{35}.\]Therefore $AB\cdot AC = (p+q)\cdot(r+s) = 560\sqrt{14}$. So the answer is $560 + 14 = \boxed{574}$

 By asr41

 



\textbf{Solution 4 (Clean)}

This solution is directly based of @CantonMathGuy's solution.We start off with a key claim.

 \textit{ Claim. } $XB \parallel AC$ and $YC \parallel AB$.

 \textit{ Proof. }

 
                    \begin{figure}[H]
                        \centering
                        \includegraphics[width=0.5\textwidth]{Solution/AIME_II_2019_15_Solution_Figure_1.png}
                    \end{figure}
                    Let $E$ and $F$ denote the reflections of the orthocenter over points $P$ and $Q$, respectively. Since $EF \parallel XY$ and 
 \[EF = 2 PQ = XP + PQ + QY = XY,\] we have that $E X Y F$ is a rectangle. Then, since $\angle XYF = 90^\circ$ we obtain $\angle XBF = 90^\circ$ (which directly follows from $XBYF$ being cyclic); hence $\angle XBQ = \angle AQB$, or $XB \parallel AQ \Rightarrow XB \parallel AC$. 

 Similarly, we can obtain $YC \parallel AB$. $\ \blacksquare$

 A direct result of this claim is that $\triangle BPX \sim \triangle APQ \sim \triangle CYQ$.

 Thus, we can set $AP = 5k$ and $BP = 2k$, then applying Power of a Point on $P$ we get $10 \cdot 40 = 10k^2 \implies k = 2\sqrt{10} \implies AB = 14 \sqrt{10}$. Also, we can set $AQ = 5l$ and $CQ = 3l$ and once again applying Power of a Point (but this time to $Q$) we get 

 $\phantom{...................}15 \cdot 35 = 15l^2 \implies l = \sqrt{35} \implies AC = 8 \sqrt{35}$. 

 Hence, 

 $\phantom{...................}AB \cdot AC = 112 \sqrt{350} = 112 \cdot 5 \sqrt{14} = 560 \sqrt{14}$ 

 and the answer is $560 + 14 = \boxed{574}$. ~rocketsri

 



\textbf{Solution 5}


                    \begin{figure}[H]
                        \centering
                        \includegraphics[width=0.5\textwidth]{Solution/AIME_II_2019_15_Solution_Figure_2.png}
                    \end{figure}
                    \textbf{mathboy282}

 



\textbf{Solution 6}

Certainly $\triangle AQP\sim\triangle ABC$. Let $k=AP$ and $\ell=AQ$, and let $c$ be the constant with $BC = PQ\cdot c = 25c$. Then $PB = \ell c-k$ and $QC = kc-\ell$. By Power of a Point at $P$ and $Q$, we obtain the equations\begin{align*}AP\cdot PB &= k(\ell c-k) = k\ell c-k^2 = 400,\\AQ\cdot QC &= \ell(kc-\ell) = k\ell c-\ell^2 = 525.\end{align*}
This implies that $k^2-\ell^2 = 125$. Next, by the Pythagorean theorem we have \begin{align*}PC^2 &= BC^2-BP^2 = 625c^2-(\ell c-k)^2,\\PC^2 &= AC^2-AP^2 = k^2c^2 - k^2.\end{align*}
Equating the two and simplifying, we arrive at the equation 
 \[c = \frac{2k\ell}{k^2+\ell^2-625}.\]From here, substitute $\ell = \sqrt{k^2-125}$. Then the first Power of a Point equation becomes 
 \[k^2+400=k\ell\cdot\frac{2k\ell}{2k^2-750} = \frac{k^2(k^2-125)}{k^2-375}.\]Solving, we obtain $k=\sqrt{1000}$ and $\ell = \sqrt{875}$. Then we evaluate:\begin{align*}AB\cdot AC &= (ck)(c\ell) = \frac{(k\ell)^3}{(k^2-375)^2}\\&= \frac{(\sqrt{875000})^3}{625^2} = \frac{(250\sqrt{14})^3}{625^2} \\&= \frac{(2\cdot 5^3)^3\cdot 14\sqrt{14}}{5^8}=560\sqrt{14}.\end{align*}
The answer is $560+14 = \boxed{574}.$




\textbf{Solution 7}

Let the circumcircle of $\triangle ABC$ be $\omega$. We know that

 $Pow_{\omega}(P) = AP \cdot BP = XP \cdot PY = 10(25 + 15) = 400$ and let $AP = a$ to get $BP = \frac{400}{a}$. We also know that

 $Pow_{\omega}(Q) = AQ \cdot CQ = XQ \cdot QY = (10 + 25)(15) = 525$ and let $AQ = b$ to get $CQ = \frac{525}{b}$.

 Let $\angle BAC = \theta$. We note that by Law of Cosines on $\triangle APQ$ we have

 $AP^2 + AQ^2 - 2\cdot AP \cdot AQ \cdot \cos(\angle BAC) = PQ^2 \implies a^2 + b^2 - 2ab\cos(\theta) = 25^2 = 625$.

 Also, we know that on $\triangle APC$
$\cos(\theta) = \frac{AP}{AC} = \frac{a}{AQ + QC} = \frac{a}{b + \frac{525}{b}} = \frac{ab}{b^2 + 525}$.

 And on $\triangle ABQ$ we know that

 $\cos(\theta) = \frac{AQ}{AB} = \frac{b}{a + \frac{400}{a}} = \frac{ab}{a^2 + 400}$.

 From here, we set the two expressions of $\cos(\theta)$ equal to immediately that that $a^2 + 400 = b^2 + 525 \implies b^2 = a^2 - 125$Also, we have

 $a^2 + b^2 - 2ab\cos(\theta) = 625 \implies a^2 + a^2 - 125 - 2ab\cos(\theta) = 625$ but we know that $\cos(\theta) = \frac{ab}{a^2 + 400}$ so $ab = (a^2 + 400)\cos(\theta)$ thus we get

 $2a^2 - 125 - 2(a^2 + 400)\cos^2(\theta) = 625 \implies a^2 - (a^2 + 400)\cos^2(\theta) = 375$ but recall that we have

 $\cos^2(\theta) = \frac{a^2 b^2}{(a^2 + 400)^2}$ so

 $a^2 - (a^2 + 400)(\frac{a^2 b^2}{(a^2 + 400)^2}) = 375 \implies a^2 - \frac{a^2 b^2}{a^2 + 400} = 375$ so we have

 $a^2 - \frac{a^2(a^2 - 125)}{a^2 + 400} = 375 \implies \frac{a^4 - 125a^2}{a^2 + 400} = a^2 - 375 \implies a^4 - 125a^2 = (a^2 - 375)(a^2 + 400)$ so we have

 $a^4 - 125a^2 = a^4 + 400a^2 - 375a^2 - 375 \cdot 400 = a^4 + 25a^2 - 375 \cdot 400$ thus the $a^4$ terms cancel out nicely and we have

 $25a^2 + 125a^2 = 375 \cdot 400 \implies a^2 = \frac{375 \cdot 400}{150} = \frac{25 \cdot 15 \cdot 40 \cdot 10}{15 \cdot 10} = 25 \cdot 40 = 100 \cdot 10$ thus $a = 10\sqrt{10}$. Recall that $b^2 = a^2 - 125 = 25 \cdot 40 - 25 \cdot 5 = 25 \cdot 35$ so $b = 5\sqrt{35}$. Finally, we have

 $AB \cdot AC = (a + \frac{400}{a})(b + \frac{525}{b}) = (\frac{a^2 + 400}{a})(\frac{b^2 + 525}{b}) = \frac{(1000 + 400)(25 \cdot 35 + 525)}{10\sqrt{10} \cdot 5\sqrt{35}}$. Note that $25 \cdot 35 = 5 \cdot 5 \cdot 5 \cdot 7 = 125 \cdot 7 = (120 + 5)(7) = 840 + 35 = 875$ and so $875 + 525 = 900 - 25 + 500 + 25 = 1400$ thus we have

 $\frac{1400 \cdot 1400}{50 \cdot \sqrt{10} \cdot \sqrt{35}} = \frac{28 \cdot 1400}{\sqrt{10} \cdot \sqrt{35}} = \frac{28 \cdot 1400}{\sqrt{5} \cdot \sqrt{2} \cdot \sqrt{5} \cdot \sqrt{7}} = \frac{28 \cdot 1400}{5 \cdot \sqrt{14}} = \frac{28 \cdot 280 \cdot \sqrt{14}}{14} = 560\sqrt{14}$ thus the answer is $560 + 14 = \boxed{574}$.

 


\end{enumerate}